\chapter{Eksperymenty numeryczne do rozdziałów 2--32}
\label{chap:pde-code}

Każdy poniższy listing jest samodzielnym programem w języku Python,
korzystającym z bibliotek \texttt{numpy} i \texttt{plotly}.
Wizualizacje są interaktywne i zapisują się do plików \texttt{.html};
pliki zawierają również lokalną kopię biblioteki JavaScript Plotly,
więc otwierają się bez połączenia z internetem.
Programy uruchamia się przez \texttt{python nazwa\_pliku.py}.
W załączonym projekcie wszystkie 29 listingów zapisano także jako
osobne pliki w katalogu \texttt{scripts/listings/}. Aby uruchomić
przykład, należy wykonać odpowiadający mu plik \texttt{.py};
kompletną tabelę nazw umieszczono w pliku
\texttt{INSTRUKCJA\_URUCHOMIENIA.md}. Eksperymenty korzystają
z NumPy w wersji co najmniej 2.0 oraz Plotly. Część przykładów,
w szczególności wizualizacja 3D wiru Taylora--Greena, wymaga ponadto
SciPy i Matplotlib. Wykonanie programu testuje jego działanie
numeryczne, ale samo nie stanowi dowodu zbieżności schematu.

\section{Charakterystyki, równanie Poissona i ciepło}
W kodzie wykorzystujemy rozwiązanie charakterystyk
równania Burgersa dla danych $-\tanh x$, tłumienie
harmonicznych modów przez $r^k$ oraz jawny profil Gaussa.

\begin{lstlisting}[style=pythonpde,caption={Burgers, funkcje harmoniczne i dyfuzja: Plotly},label={code:pde-old}]
"""Interaktywne wykresy PDE: charakterystyki, profile, dyfuzja i Poisson.
Wymagania: numpy, plotly. Zapisuje samodzielne pliki HTML (offline).
"""
from pathlib import Path
import numpy as np
import plotly.graph_objects as go

OUT = Path(__file__).resolve().parent


def save(fig: go.Figure, name: str) -> None:
    fig.update_layout(template="plotly_white")
    fig.write_html(OUT / name, include_plotlyjs=True, full_html=True)


def data(x):
    return -np.tanh(x)


def position(t, xi):
    return xi + t * data(xi)


def burgers(t, x):
    xi = np.linspace(-8.0, 8.0, 100001)
    X = position(t, xi)
    if not np.all(np.diff(X) > 0):
        raise ValueError("Brak rozwiazania klasycznego dla tego czasu")
    return data(np.interp(x, X, xi))


def characteristic_plot():
    fig = go.Figure()
    ts = np.linspace(0, 0.99, 250)
    for xi in np.linspace(-2.4, 2.4, 13):
        fig.add_trace(go.Scatter(x=position(ts, xi), y=ts,
                                 mode="lines", name=f"xi={xi:.1f}"))
    fig.update_layout(title="Charakterystyki Burgersa", xaxis_title="x",
                      yaxis_title="t")
    save(fig, "burgers_characteristics.html")


def profile_plot():
    x = np.linspace(-3.0, 3.0, 1001)
    fig = go.Figure()
    for t in (0.0, 0.5, 0.9, 0.99):
        fig.add_trace(go.Scatter(x=x, y=burgers(t, x),
                                 mode="lines", name=f"t={t:.2f}"))
    fig.update_layout(title="Burgers: profile przed utrata regularnosci",
                      xaxis_title="x", yaxis_title="u(t,x)")
    save(fig, "burgers_profiles.html")


def poisson_modes():
    k = np.arange(0, 30)
    fig = go.Figure()
    for r in (0.2, 0.5, 0.8):
        fig.add_trace(go.Scatter(x=k, y=r**k, mode="lines+markers",
                                 name=f"r={r}"))
    fig.update_layout(title="Wygaszanie modow harmonicznych",
                      xaxis_title="czestosc k", yaxis_title="amplituda r^k",
                      yaxis_type="log")
    save(fig, "poisson_modes.html")


def heat_gaussian_plot():
    x = np.linspace(-4, 4, 1200)
    fig = go.Figure()
    for t in (0.0, 0.10, 0.40, 1.20):
        u = np.exp(-x**2 / (1 + 4*t)) / np.sqrt(1 + 4*t)
        fig.add_trace(go.Scatter(x=x, y=u, mode="lines", name=f"t={t:.2f}"))
    fig.update_layout(title="Dokladne rozwiazanie Gaussowskie (kappa=1)",
                      xaxis_title="x", yaxis_title="u(t,x)")
    save(fig, "heat_gaussian_profiles.html")


if __name__ == "__main__":
    characteristic_plot()
    profile_plot()
    poisson_modes()
    heat_gaussian_plot()
    for t in (0.0, 0.5, 0.9, 0.99):
        print(f"t={t:.2f}, pochodna w x=0: {-1/(1-t):.6f}")
\end{lstlisting}

\section{Schemat falowy i kontrola dokładności}
\begin{lstlisting}[style=pythonpde,caption={Schemat różnicowy równania falowego: Plotly i test błędu},label={code:pde-wave}]
"""Jawny schemat falowy: interaktywny wykres Plotly i blad max.
Wymagania: numpy, plotly. Zapisuje plik HTML do katalogu skryptu.
"""
from pathlib import Path
import numpy as np
import plotly.graph_objects as go

c, L, N = 1.0, 1.0, 300
dx = L / N
lam = 0.8
dt = lam * dx / c
x = np.linspace(0.0, L, N + 1)
T = 0.8
steps = int(round(T / dt))


def exact(t):
    return (np.cos(2*np.pi*c*t/L)*np.sin(2*np.pi*x/L)
            + 0.4*np.cos(7*np.pi*c*t/L)*np.sin(7*np.pi*x/L))


u_prev = exact(0.0)
u = u_prev.copy()
u[1:-1] += 0.5*lam**2*(u_prev[2:] - 2*u_prev[1:-1] + u_prev[:-2])
u[0] = u[-1] = 0.0
snapshots = {0: u_prev.copy()}
for m in range(1, steps):
    u_next = np.zeros_like(u)
    u_next[1:-1] = (2*u[1:-1] - u_prev[1:-1]
                    + lam**2*(u[2:] - 2*u[1:-1] + u[:-2]))
    u_prev, u = u, u_next
    if m + 1 in (steps//3, 2*steps//3, steps):
        snapshots[m + 1] = u.copy()

fig = go.Figure()
for m, values in snapshots.items():
    fig.add_trace(go.Scatter(x=x, y=values, mode="lines", name=f"t={m*dt:.3f}"))
fig.update_layout(template="plotly_white", title="Struna: lambda=0.8",
                  xaxis_title="x", yaxis_title="u(t,x)")
fig.write_html(Path(__file__).with_name("wave_leapfrog.html"),
               include_plotlyjs=True, full_html=True)
print("CFL:", lam, "blad max:", np.max(np.abs(u - exact(steps*dt))))
\end{lstlisting}

\section{Szeregi Fouriera i ewolucje modalne}
\begin{lstlisting}[style=pythonpde,caption={Średnie Fejéra, Gibbs oraz powierzchnie ewolucji: Plotly},label={code:pde-fourier}]
"""Szeregi Fouriera i PDE na odcinku, wykresy interaktywne Plotly.
Wymagania: numpy, plotly. Zapisuje samodzielne pliki HTML.
"""
from pathlib import Path
import numpy as np
import plotly.graph_objects as go

OUT = Path(__file__).resolve().parent


def save(fig, name):
    fig.update_layout(template="plotly_white")
    fig.write_html(OUT / name, include_plotlyjs=True, full_html=True)


def fourier_vs_fejer():
    # f(x) = sign(x), x in (-pi, pi). W wartosciach skoku: f=0.
    x = np.linspace(-np.pi, np.pi, 1400)
    f = np.sign(x)
    fig = go.Figure([go.Scatter(x=x, y=f, name="dane skokowe")])
    for N in (5, 15, 45):
        ks = np.arange(1, N+1, 2)
        terms = (4/np.pi) * np.sin(np.outer(ks, x))/ks[:, None]
        fourier = terms.sum(axis=0)
        fejer = ((1-ks/(N+1))[:, None] * terms).sum(axis=0)
        fig.add_trace(go.Scatter(x=x, y=fourier, name=f"S_{N}"))
        fig.add_trace(go.Scatter(x=x, y=fejer, name=f"sigma_{N}"))
    fig.update_layout(title="Zjawisko Gibbsa i srednie Fejera",
                      xaxis_title="x", yaxis_title="wartosc")
    save(fig, "fourier_vs_fejer.html")


def heat_modes():
    L, kappa = 1.0, 0.06
    x = np.linspace(0.0, L, 500)
    t = np.linspace(0.0, 0.6, 100)
    X, T = np.meshgrid(x, t)
    U = (np.exp(-kappa*np.pi**2*T)*np.sin(np.pi*X)
         + 0.45*np.exp(-kappa*(5*np.pi)**2*T)*np.sin(5*np.pi*X))
    fig = go.Figure(data=go.Surface(x=X, y=T, z=U, colorscale="Viridis"))
    fig.update_layout(title="Rownanie ciepla: tlumienie wysokich modow",
                      scene=dict(xaxis_title="x", yaxis_title="t",
                                 zaxis_title="u(t,x)"))
    save(fig, "heat_modes.html")


def wave_modes():
    L, c = 1.0, 1.0
    x = np.linspace(0.0, L, 500)
    t = np.linspace(0.0, 2.0, 140)
    X, T = np.meshgrid(x, t)
    U = (np.cos(np.pi*c*T/L)*np.sin(np.pi*X/L)
         + 0.45*np.cos(5*np.pi*c*T/L)*np.sin(5*np.pi*X/L))
    fig = go.Figure(data=go.Surface(x=X, y=T, z=U, colorscale="Viridis"))
    fig.update_layout(title="Rownanie falowe: utrzymywanie sie modow",
                      scene=dict(xaxis_title="x", yaxis_title="t",
                                 zaxis_title="u(t,x)"))
    save(fig, "wave_modes.html")


if __name__ == "__main__":
    fourier_vs_fejer()
    heat_modes()
    wave_modes()
    print("Wykresy Plotly zapisane w", OUT)
\end{lstlisting}

\section{Przybliżone jedynki, skok i jądro cieplne}
\begin{lstlisting}[style=pythonpde,caption={Dystrybucje: przybliżone jedynki, funkcja Heaviside'a i jądro cieplne (Plotly)},label={code:pde-distributions}]
"""Przyblizone jedynki, pochodne slabego skoku i jadro ciepla.
Wymagania: numpy, plotly. Eksportuje interaktywne HTML.
"""
from pathlib import Path
import numpy as np
import plotly.graph_objects as go

OUT = Path(__file__).resolve().parent
x = np.linspace(-2.0, 2.0, 10001)
dx = x[1] - x[0]


def bump(z):
    """Gladki profil o zwartym nosniku [-1,1], bez normalizacji."""
    z = np.asarray(z)
    y = np.zeros_like(z, dtype=float)
    mask = np.abs(z) < 1.0
    y[mask] = np.exp(-1.0 / (1.0 - z[mask] ** 2))
    return y


zq = np.linspace(-1.0, 1.0, 5001)
MASS = np.trapezoid(bump(zq), zq)


def mollifier(epsilon):
    return bump(x / epsilon) / (epsilon * MASS)


def save(fig, name):
    fig.update_layout(template="plotly_white", hovermode="x unified")
    fig.write_html(OUT / name, include_plotlyjs=True, full_html=True)


def delta_and_heaviside():
    fig_delta, fig_h = go.Figure(), go.Figure()
    eps_values = (0.5, 0.25, 0.10)
    for eps in eps_values:
        rho = mollifier(eps)
        fig_delta.add_trace(go.Scatter(x=x, y=rho, mode="lines",
                                       name=f"epsilon={eps}"))
        # H * rho = integral od -nieskonczonosci do x z rho(s) ds
        h = np.r_[0.0, np.cumsum((rho[:-1] + rho[1:]) * dx / 2)]
        fig_h.add_trace(go.Scatter(x=x, y=h, mode="lines",
                                   name=f"epsilon={eps}"))
        # Test dystrybucyjny: int rho_epsilon(x) * phi(x) dx -> phi(0)=1
        phi = np.exp(-x*x)
        pairing = np.trapezoid(rho * phi, x)
        print(f"epsilon={eps:5.2f}, masa={np.trapezoid(rho,x):.7f}, "
              f"<rho_epsilon,phi>={pairing:.7f}, phi(0)=1")
    fig_h.add_trace(go.Scatter(x=x, y=(x > 0).astype(float),
                               mode="lines", name="H (skok)",
                               line=dict(dash="dot", color="gray")))
    fig_delta.update_layout(title="Gladkie aproksymacje delty Diraca",
                            xaxis_title="x", yaxis_title="rho_epsilon(x)")
    fig_h.update_layout(title="Wygladzona funkcja Heaviside'a",
                        xaxis_title="x", yaxis_title="(H * rho_epsilon)(x)")
    save(fig_delta, "pde7_delta_plotly.html")
    save(fig_h, "pde7_heaviside_plotly.html")


def heat_kernel():
    kappa = 0.15
    fig = go.Figure()
    for t in (0.01, 0.04, 0.2):
        gamma = np.exp(-x*x/(4*kappa*t)) / np.sqrt(4*np.pi*kappa*t)
        fig.add_trace(go.Scatter(x=x, y=gamma, name=f"t={t}"))
        print(f"t={t:0.3f}, masa na [-2,2]={np.trapezoid(gamma,x):.7f}")
    fig.update_layout(title="Jadro cieplne w 1D jako aproksymacja delty",
                      xaxis_title="x", yaxis_title="Gamma(t,x)")
    save(fig, "pde7_heat_kernel_plotly.html")


if __name__ == "__main__":
    delta_and_heaviside()
    heat_kernel()
\end{lstlisting}


\section[Transformata, Parseval i mnożniki]{Transformata na $\R^d$, tożsamość Parsevala i mnożniki}
\begin{lstlisting}[style=pythonpde,
caption={Transformata Gaussa, dyskretna tożsamość Parsevala, ewolucje spektralne oraz normy $H^s$ (Plotly)},
label={code:pde8-plotly}]
"""PDE, rozdzial 8: FFT, Plancherel, ewolucja ciepla i fal, Plotly.
Wymagania: numpy>=2.0, plotly>=5. Zapisuje samodzielne pliki HTML.
Transformata ciagla: F f(xi)=int exp(-i x xi) f(x) dx.
Siatka FFT jest okresowa: L=N*dx, xi=2*pi*fftfreq(N,d=dx).
"""
from pathlib import Path
import numpy as np
import plotly.graph_objects as go
OUT = Path(__file__).resolve().parent
OUT.mkdir(parents=True, exist_ok=True)

def save_plot(fig, filename):
    fig.update_layout(template="plotly_white", hovermode="x unified",
                      margin=dict(l=60, r=25, t=65, b=60))
    fig.write_html(OUT / filename, include_plotlyjs=True, full_html=True)

def gauss_fft():
    n, length, sigma = 2048, 24.0, 1.1
    dx = length / n
    x = (np.arange(n) - n//2) * dx
    f = np.exp(-x*x/(2*sigma*sigma))
    xi = np.fft.fftshift(np.fft.fftfreq(n, d=dx))*2*np.pi
    fhat = dx * np.fft.fftshift(np.fft.fft(np.fft.ifftshift(f)))
    exact = np.sqrt(2*np.pi)*sigma*np.exp(-sigma*sigma*xi*xi/2)
    mask = np.abs(xi) <= 8.0
    err = np.max(np.abs(fhat[mask]-exact[mask]))
    parseval_x = dx * np.sum(np.abs(f)**2)
    dxi = 2*np.pi/length
    parseval_k = (dxi/(2*np.pi))*np.sum(np.abs(fhat)**2)
    print("FFT vs Gauss: blad max na |xi|<=8:", f"{err:.3e}")
    print("Parseval: przestrzen", f"{parseval_x:.10f}",
          "czestotliwosc", f"{parseval_k:.10f}")
    fig = go.Figure()
    fig.add_trace(go.Scatter(x=xi[mask], y=exact[mask], mode="lines",
                             name="transformata dokladna"))
    fig.add_trace(go.Scatter(x=xi[mask], y=np.real(fhat[mask]), mode="markers",
                             marker=dict(size=4), name="FFT * dx"))
    fig.update_layout(title="Transformata Gaussa: wzor dokladny i FFT",
                      xaxis_title="czestotliwosc xi",
                      yaxis_title="F[f](xi)")
    save_plot(fig, "pde8_gauss_fft.html")

def modes_heat_wave():
    n, length = 1024, 2*np.pi
    dx = length/n
    x = np.arange(n)*dx
    frequencies = 2*np.pi*np.fft.fftfreq(n, d=dx)
    # Dane na torusie: te same dla obu rownan, predkosc poczatkowa fali 0.
    f = np.sin(2*x) + 0.55*np.cos(9*x) + 0.25*np.sin(23*x)
    fhat = np.fft.fft(f)
    kappa, c = 0.09, 1.0
    times = (0.0, 0.08, 0.3, 0.85)
    fig_heat, fig_wave = go.Figure(), go.Figure()
    for t in times:
        heat = np.fft.ifft(np.exp(-kappa*t*frequencies**2)*fhat).real
        wave = np.fft.ifft(np.cos(c*t*np.abs(frequencies))*fhat).real
        fig_heat.add_trace(go.Scatter(x=x, y=heat, name=f"t={t:g}"))
        fig_wave.add_trace(go.Scatter(x=x, y=wave, name=f"t={t:g}"))
    fig_heat.update_layout(title="Rownanie ciepla: tlumienie wysokich modow",
                           xaxis_title="x", yaxis_title="u(t,x)")
    fig_wave.update_layout(title="Rownanie falowe: mody oscyluja",
                           xaxis_title="x", yaxis_title="u(t,x)")
    save_plot(fig_heat, "pde8_heat_fft.html")
    save_plot(fig_wave, "pde8_wave_fft.html")

def h_s_delta():
    # Obciecie czestotliwosci: dla s<-1/2 calka stabilizuje sie.
    s_values = (-0.8, -0.5, -0.2)
    cutoffs = np.geomspace(2, 250, 60)
    # Dla d=1, ||delta||_{H^s}^2=(2*pi)^-1 * int(1+xi^2)^s dxi.
    xi = np.linspace(0, 250, 20001)
    fig = go.Figure()
    for s in s_values:
        vals = (1+xi*xi)**s
        increments = 0.5*(vals[1:]+vals[:-1])*np.diff(xi)
        cumulative = np.r_[0.0, np.cumsum(increments)] / np.pi
        approx = np.interp(cutoffs, xi, cumulative)
        fig.add_trace(go.Scatter(x=cutoffs, y=approx, name=f"s={s}"))
    fig.update_layout(title="Delta w H^s(R): normy z obcieciem czestotliwosci",
                      xaxis_title="odciecie |xi| <= R", xaxis_type="log",
                      yaxis_title="(2*pi)^-1 calka do R")
    save_plot(fig, "pde8_delta_Hs.html")
if __name__ == "__main__":
    gauss_fft()
    modes_heat_wave()
    h_s_delta()
    print("Wyniki zapisane w", OUT)
\end{lstlisting}

\section{Wygładzanie Sobolewa i warunki brzegowe}
Poniższy eksperyment porównuje zbieżność splotu w normie $H^1$
ze zbieżnością obcięć brzegowych jedynie w $L^2$.
\begin{lstlisting}[style=pythonpde,caption={Wygładzanie w $H^1$, pochodna słaba i obcięcia brzegowe -- Plotly},label={code:pde-sobolev9}]
"""Iteracja 9: Sobolev W^{1,2}, periodyczny splot i warunki brzegowe.

Wymagania: numpy, plotly. Zapisuje trzy samodzielne pliki HTML.
"""
from pathlib import Path
import numpy as np
import plotly.graph_objects as go

HERE = Path(__file__).resolve().parent


def save(fig, filename):
    fig.update_layout(template="plotly_white", hovermode="x unified",
                      font=dict(size=13), legend=dict(orientation="h", y=1.11))
    fig.write_html(HERE / filename, include_plotlyjs=True, full_html=True)


def periodic_mollification():
    n = 4096
    x = -np.pi + 2 * np.pi * np.arange(n) / n
    h = 2 * np.pi / n
    u = np.abs(np.sin(x))
    du = np.sign(np.sin(x)) * np.cos(x)  # pochodna slaba a.e.
    k = np.fft.fftfreq(n, d=h) * (2 * np.pi)
    uhat = np.fft.fft(u)
    epsilons = [0.40, 0.20, 0.10, 0.05, 0.025, 0.0125]
    fig = go.Figure()
    fig.add_trace(go.Scatter(x=x, y=u, name="u=|sin(x)|", line=dict(width=3)))
    errors = []
    for eps in epsilons:
        mult = np.exp(-0.5 * eps**2 * k**2)
        smooth = np.fft.ifft(mult * uhat).real
        d_smooth = np.fft.ifft(1j * k * mult * uhat).real
        err_l2 = np.sqrt(h * np.sum((smooth - u)**2))
        err_grad = np.sqrt(h * np.sum((d_smooth - du)**2))
        errors.append((eps, np.hypot(err_l2, err_grad)))
        if eps in [0.40, 0.10, 0.025]:
            fig.add_trace(go.Scatter(x=x, y=smooth, name=f"eps={eps:.3f}"))
    fig.update_layout(title="Wygladzanie w H^1(T): u(x)=|sin x|",
                      xaxis_title="x", yaxis_title="u_epsilon(x)")
    save(fig, "pde9_wygladzanie_plotly.html")

    fig2 = go.Figure()
    fig2.add_trace(go.Scatter(x=x, y=du, name="pochodna slaba", line=dict(width=2)))
    for eps in [0.40, 0.10, 0.025]:
        derivative = np.fft.ifft(1j * k * np.exp(-0.5*eps**2*k**2) * uhat).real
        fig2.add_trace(go.Scatter(x=x, y=derivative, name=f"eps={eps:.3f}"))
    fig2.update_layout(title="Wygladzanie pochodnej slabej",
                       xaxis_title="x", yaxis_title="d u_epsilon / dx")
    save(fig2, "pde9_pochodne_plotly.html")
    print("Spadek bledu H1:", errors)
    assert errors[-1][1] < errors[0][1], "Blad powinien malec przy eps->0"


def smoothstep(z):
    v = np.clip(z, 0.0, 1.0)
    return 3*v*v - 2*v*v*v


def smoothstep_deriv(z):
    v = np.clip(z, 0.0, 1.0)
    return 6*v*(1-v)


def boundary_cutoffs():
    x = (np.arange(20001) + 0.5) / 20001
    vals = []
    fig = go.Figure()
    for n in [4, 8, 16, 32, 64]:
        a, b = smoothstep(n*x), smoothstep(n*(1-x))
        d_a = n*smoothstep_deriv(n*x)
        d_b = -n*smoothstep_deriv(n*(1-x))
        chi = a*b
        d_chi = d_a*b + a*d_b
        err = np.sqrt(np.mean((chi-1)**2))
        grad = np.sqrt(np.mean(d_chi**2))
        vals.append((n, err, grad))
        fig.add_trace(go.Scatter(x=x, y=chi, name=f"n={n}"))
    fig.update_layout(title="Obciecia daza do 1 w L^2, ale nie w H^1",
                      xaxis_title="x w (0,1)", yaxis_title="chi_n(x)")
    save(fig, "pde9_slady_obciecia_plotly.html")
    print("n, ||chi-1||L2, ||chi'||L2:", vals)
    assert vals[-1][1] < vals[0][1]
    assert vals[-1][2] > vals[0][2]


if __name__ == "__main__":
    periodic_mollification()
    boundary_cutoffs()
\end{lstlisting}


\section{Koncentracja krytyczna i ciąg Mosera — Plotly}
Poniższy program wizualizuje koncentrację przy wykładniku $p^*=6$ dla $d=3$, $p=2$, porównuje skalowanie norm oraz ciąg Mosera w wymiarze $2$. Wyniki są zapisywane w trzech interaktywnych plikach HTML.
\begin{lstlisting}[style=pythonpde,caption={Koncentracja i ostrosc zanurzen Sobolewa - Plotly},label={code:pde-embed10}]
"""PDE, rozdzial 10. Interaktywne wizualizacje zanurzen Sobolewa.
Uruchom: python demo_zanurzenia_plotly.py
Pliki HTML zapisuja lokalna kopie Plotly.js i dzialaja offline.
"""
from pathlib import Path
import numpy as np
import plotly.graph_objects as go

OUT = Path(__file__).resolve().parent
EPS = [0.8, 0.4, 0.2, 0.1, 0.05]


def bump(r):
    r = np.asarray(r, dtype=float)
    v = np.zeros_like(r)
    idx = np.abs(r) < 1
    v[idx] = np.exp(-1 / (1 - r[idx] ** 2))
    return v


def bump_prime(r):
    r = np.asarray(r, dtype=float)
    v = np.zeros_like(r)
    idx = np.abs(r) < 1
    v[idx] = bump(r[idx]) * (-2 * r[idx] / (1 - r[idx] ** 2) ** 2)
    return v


# d=3, p=2. Wykladnik krytyczny p*=6, skala eps^{-1/2}.
x = np.linspace(-1, 1, 1601)
fig = go.Figure()
for eps in EPS:
    fig.add_trace(go.Scatter(
        x=x, y=eps ** (-0.5) * bump(x / eps),
        name=f"eps={eps:g}", mode="lines"))
fig.update_layout(title="Koncentracja: W1,2 w wymiarze 3 (przekroj radialny)",
                  xaxis_title="x", yaxis_title="u_eps(x)", template="plotly_white")
fig.write_html(OUT / "pde10_koncentracja.html", include_plotlyjs=True)

# Calkowanie sferyczne: 4*pi int_0^1 |v(r)|^q r2 dr.
r = np.linspace(0, 1, 50001)
q_values = [2, 4, 6, 8]
baseline = {q: (4 * np.pi * np.trapezoid(bump(r) ** q * r ** 2, r)) ** (1 / q)
            for q in q_values}
grad = np.sqrt(4 * np.pi * np.trapezoid(bump_prime(r) ** 2 * r ** 2, r))
fig = go.Figure()
for q in q_values:
    vals = [e ** (-0.5 + 3 / q) * baseline[q] for e in EPS]
    fig.add_trace(go.Scatter(x=EPS, y=vals, mode="lines+markers", name=f"L^{q}"))
fig.add_trace(go.Scatter(x=EPS, y=[grad] * len(EPS),
                         mode="lines+markers", name="norma gradientu L2"))
fig.update_layout(title="Skalowanie norm: krytyczne L6 nie zalezy od eps",
                  xaxis_title="eps", yaxis_title="norma",
                  xaxis_type="log", yaxis_type="log", template="plotly_white")
fig.write_html(OUT / "pde10_skalowanie_norm.html", include_plotlyjs=True)

# Ciag Mosera w wymiarze 2: sup rosnie, norma gradientu stala.
r = np.geomspace(1e-5, 1, 900)
fig = go.Figure()
for eps in [0.1, 0.01, 0.001, 0.0001]:
    L = np.log(1 / eps)
    val = np.minimum(np.log(1 / r), L) / np.sqrt(L)
    fig.add_trace(go.Scatter(x=r, y=val, mode="lines", name=f"eps={eps:g}"))
fig.update_layout(title="Ciag Mosera: wzrost maksimum przy stalej normie gradu w L2",
                  xaxis_title="promien r", yaxis_title="u_eps(r)",
                  xaxis_type="log", template="plotly_white")
fig.write_html(OUT / "pde10_moser.html", include_plotlyjs=True)

for q in q_values:
    assert np.isfinite(baseline[q]) and baseline[q] > 0
assert np.isfinite(grad) and grad > 0
print(f"Kontrola norm: L6(eps=0.8)={baseline[6]:.8f}; "
      f"L6(eps=0.05)={baseline[6]:.8f}; grad L2={grad:.8f}")
print("Wygenerowano: pde10_koncentracja.html, pde10_skalowanie_norm.html, pde10_moser.html")
\end{lstlisting}


\section{Galerkin, zbieżność i energia — Plotly}
\label{sec:pde11-code}
Poniższy program tworzy macierz sztywności metody elementów liniowych,
porównuje aproksymację z rozwiązaniem dokładnym, weryfikuje wzory
na błędy w $L^2$ i seminormie $H^1$ oraz wizualizuje funkcjonał energii.
Wynikami są trzy interaktywne pliki HTML, otwierane bez połączenia
z Internetem.
\begin{lstlisting}[style=pythonpde,caption={Galerkin, błędy i funkcjonał energii (Plotly)},label={code:pde11-galerkin}]
"""PDE - rozdzial 11: metoda Galerkina dla -u''=1, u(0)=u(1)=0.

Instalacja: pip install numpy plotly
Uruchom: python demo_wariacyjne_plotly.py
Wyniki: trzy samodzielne pliki HTML z Plotly (offline).
"""
from pathlib import Path

import numpy as np
import plotly.graph_objects as go

OUT = Path(__file__).resolve().parent


def exact(x):
    return x * (1.0 - x) / 2.0


def fem_solution(N):
    """Zwraca wezly i wartosci liniowych elementow skonczonych."""
    h = 1.0 / N
    diagonal = np.full(N - 1, 2.0 / h)
    K = np.diag(diagonal)
    if N > 2:
        offdiag = np.full(N - 2, -1.0 / h)
        K += np.diag(offdiag, 1) + np.diag(offdiag, -1)
    b = np.full(N - 1, h)
    U = np.linalg.solve(K, b)
    nodes = np.linspace(0.0, 1.0, N + 1)
    values = np.r_[0.0, U, 0.0]
    assert np.allclose(values, exact(nodes), atol=1e-12)
    assert np.linalg.eigvalsh(K)[0] > 0
    return nodes, values


def quadrature_errors(nodes, values):
    """Bledy L2 i H1 liczone calkowaniem Gaussa na kazdej komorce."""
    z, weights = np.polynomial.legendre.leggauss(4)
    sq_l2 = sq_h1 = 0.0
    for i in range(len(nodes) - 1):
        x0, x1 = nodes[i : i + 2]
        h = x1 - x0
        x = (x0 + x1) / 2.0 + h * z / 2.0
        slope = (values[i + 1] - values[i]) / h
        approx = values[i] + slope * (x - x0)
        sq_l2 += h / 2.0 * np.sum(weights * (exact(x) - approx) ** 2)
        sq_h1 += h / 2.0 * np.sum(weights * ((0.5 - x) - slope) ** 2)
    return np.sqrt(sq_l2), np.sqrt(sq_h1)


def save(fig, name):
    fig.update_layout(template="plotly_white", font=dict(size=13))
    fig.write_html(OUT / name, include_plotlyjs=True, full_html=True)


def main():
    Ns = [4, 8, 16, 32, 64]
    h_values, errors_l2, errors_h1 = [], [], []
    fine = np.linspace(0.0, 1.0, 1200)

    profile = go.Figure()
    profile.add_trace(go.Scatter(x=fine, y=exact(fine), name="Rozwiazanie dokladne"))

    for N in Ns:
        nodes, values = fem_solution(N)
        err_l2, err_h1 = quadrature_errors(nodes, values)
        h = 1.0 / N
        expected_l2 = h**2 / np.sqrt(120.0)
        expected_h1 = h / np.sqrt(12.0)
        assert np.isclose(err_l2, expected_l2, rtol=1e-9, atol=1e-12)
        assert np.isclose(err_h1, expected_h1, rtol=1e-9, atol=1e-12)
        h_values.append(h)
        errors_l2.append(err_l2)
        errors_h1.append(err_h1)
        if N <= 16:
            profile.add_trace(go.Scatter(x=nodes, y=values, name=f"FEM N={N}",
                                         mode="lines+markers"))
        print(f"N={N:2d}, h={h:.5f}, blad L2={err_l2:.8e}, blad H1={err_h1:.8e}")

    profile.update_layout(title="Galerkin: rozwiazanie dokladne i aproksymacje liniowe",
                          xaxis_title="x", yaxis_title="u(x)")
    save(profile, "pde11_fem_rozwiazania.html")

    convergence = go.Figure()
    convergence.add_trace(go.Scatter(x=h_values, y=errors_l2,
                                     name="||u-u_h||_L2 (rzad 2)", mode="lines+markers"))
    convergence.add_trace(go.Scatter(x=h_values, y=errors_h1,
                                     name="||u'-u_h'||_L2 (rzad 1)", mode="lines+markers"))
    convergence.update_layout(title="Zbieznosc FEM: dokladne bledy dla -u''=1",
                              xaxis_title="h", yaxis_title="norma bledu",
                              xaxis_type="log", yaxis_type="log")
    save(convergence, "pde11_fem_zbieznosc.html")

    # Dla u(x)=x(1-x)/2 mamy int_0^1 |u'|^2 dx = int_0^1 u dx = 1/12.
    t = np.linspace(-0.25, 2.25, 300)
    energy = (0.5 * t * t - t) / 12.0
    energy_plot = go.Figure()
    energy_plot.add_trace(go.Scatter(x=t, y=energy, name="J(tu)"))
    energy_plot.add_trace(go.Scatter(x=[1], y=[-1.0 / 24.0], mode="markers",
                                     marker=dict(size=11), name="Minimum t=1"))
    energy_plot.update_layout(title="Zasada Dirichleta: energia na prostej tu",
                              xaxis_title="t", yaxis_title="J(tu)")
    save(energy_plot, "pde11_energia.html")

    orders = np.polyfit(np.log(h_values), np.log(errors_l2), 1)[0], \
             np.polyfit(np.log(h_values), np.log(errors_h1), 1)[0]
    assert np.allclose(orders, (2.0, 1.0), atol=1e-9)
    print(f"Zweryfikowane rzedy: L2={orders[0]:.12f}, H1={orders[1]:.12f}")


if __name__ == "__main__":
    main()
\end{lstlisting}

\section{Spektrum, półgrupy i rezonans — Plotly}
\label{sec:pde12-code}
Poniższy program porównuje dokładne wartości własne operatora
$-\frac{\mathrm d^2}{\mathrm dx^2}$ na $(0,1)$ z widmem macierzy
różnicowej, wizualizuje ewolucję cieplną i falową tych samych
danych początkowych oraz wzrost normy rezolwenty w pobliżu
pierwszej wartości własnej. Sprawdza również rząd zbieżności
wartości własnych, równy $2$ dla zastosowanego schematu.
\begin{lstlisting}[style=pythonpde,caption={Widmo operatora Dirichleta, ewolucje i rezonans (Plotly)},label={code:pde12-operators}]
"""PDE, rozdzial 12: widmo laplasjanu Dirichleta i ewolucje.

Instalacja: pip install numpy plotly
Uruchom: python demo_operatory_plotly.py
Wyniki: trzy samodzielne pliki HTML, dzialajace offline.
"""
from pathlib import Path
import numpy as np
import plotly.graph_objects as go

OUT = Path(__file__).resolve().parent
PI = np.pi


def discrete_spectrum(n_intervals):
    h = 1.0 / n_intervals
    matrix = np.diag(np.full(n_intervals - 1, 2.0 / h**2))
    matrix += np.diag(np.full(n_intervals - 2, -1.0 / h**2), 1)
    matrix += np.diag(np.full(n_intervals - 2, -1.0 / h**2), -1)
    values = np.linalg.eigvalsh(matrix)
    j = np.arange(1, n_intervals, dtype=float)
    formula = 4.0 / h**2 * np.sin(j * PI * h / 2.0) ** 2
    assert np.allclose(values, formula, rtol=2e-12, atol=2e-11)
    return values


def save(fig, name):
    fig.update_layout(template="plotly_white", font=dict(size=13))
    fig.write_html(OUT / name, include_plotlyjs=True, full_html=True)


def main():
    sizes = [12, 24, 48, 96, 192]
    nplot = 6
    exact = (PI * np.arange(1, nplot + 1)) ** 2
    fig = go.Figure()
    fig.add_trace(go.Scatter(x=np.arange(1, nplot + 1), y=exact,
                             name="Wartosci dokladne", mode="lines+markers"))
    errors = []
    for n in sizes:
        eigen = discrete_spectrum(n)
        fig.add_trace(go.Scatter(x=np.arange(1, nplot + 1), y=eigen[:nplot],
                                 name=f"FD, N={n}", mode="lines+markers"))
        errors.append(abs(eigen[0] - PI**2))
    fig.update_layout(title="Laplasjan Dirichleta: widmo ciagle vs. macierz FD",
                      xaxis_title="Numer modu j", yaxis_title="lambda_j")
    save(fig, "pde12_widmo_plotly.html")
    h = 1.0 / np.array(sizes, dtype=float)
    order = np.polyfit(np.log(h), np.log(errors), 1)[0]
    print(f"Rzad bledu pierwszej wartosci wlasnej = {order:.4f}")
    assert 1.95 < order < 2.05

    x = np.linspace(0.0, 1.0, 1000)
    modes = np.array([1, 3, 6], dtype=float)
    coeffs = np.array([1.0, 0.35, -0.2])
    basis = np.sin(np.pi * modes[:, None] * x[None, :])
    lambdas = (np.pi * modes) ** 2
    fig2 = go.Figure()
    for t in (0.0, 0.01, 0.05, 0.12):
        y = np.sum((coeffs * np.exp(-t * lambdas))[:, None] * basis, axis=0)
        fig2.add_trace(go.Scatter(x=x, y=y, name=f"Cieplo t={t:g}"))
    for t in (0.05, 0.12):
        y = np.sum((coeffs * np.cos(t * np.sqrt(lambdas)))[:, None] * basis,
                   axis=0)
        fig2.add_trace(go.Scatter(x=x, y=y, name=f"Fala t={t:g}",
                                  line=dict(dash="dash")))
    fig2.update_layout(title="Te same dane: wygladzanie ciepla i oscylacje fali",
                       xaxis_title="x", yaxis_title="u(t,x)")
    save(fig2, "pde12_cieplo_fala_plotly.html")

    lambda1 = PI**2
    delta = np.geomspace(0.001, 4.0, 160)
    fig3 = go.Figure()
    fig3.add_trace(go.Scatter(x=delta, y=1.0 / delta,
                              name="|| (A-(lambda_1+delta)I)^-1 phi_1 ||"))
    fig3.update_layout(title="Alternatywa Fredholma: rezonans blisko lambda_1",
                       xaxis_title="delta = |zeta-lambda_1|",
                       yaxis_title="Norma rozwiazania", xaxis_type="log",
                       yaxis_type="log")
    save(fig3, "pde12_rezolwenta_plotly.html")
    print("Zweryfikowano formula FD, blad O(h^2) i trzy wizualizacje Plotly.")


if __name__ == "__main__":
    main()
\end{lstlisting}


\section{Eliptyczne zagadnienie transmisyjne — Plotly}
\label{sec:pde13-code}
Program rozwiązuje metodą elementów skończonych zadanie
$-(a u')'=1$ na $(-1,1)$, z warunkami Dirichleta i skokiem
współczynnika $a$ w zerze. Porównuje rozwiązanie numeryczne ze
wzorem~\eqref{eq:pde13-exact}, przedstawia strumień
$q=-a u'$ oraz bada błędy w normie $L^2$ i seminormie $H^1$.
Sprawdza także dokładne wzory na błędy interpolacji:
$\|u-u_h\|_{L^2}=h^2\sqrt{17/1920}$ i
$\|(u-u_h)'\|_{L^2}=h\sqrt{17/192}$.
\begin{lstlisting}[style=pythonpde,caption={Transmisja i zbieżność metody Galerkina (Plotly)},label={code:pde13-transmission}]
"""PDE, rozdzial 13: dywergencyjne zagadnienie transmisyjne.

Instalacja: pip install numpy plotly
Uruchom: python demo_eliptyczne_plotly.py
Wyniki: samodzielne pliki HTML (dzialaja offline).
"""
from pathlib import Path
import numpy as np
import plotly.graph_objects as go

OUT = Path(__file__).resolve().parent


def exact(x):
    x = np.asarray(x, dtype=float)
    return np.where(x <= 0, -x*x/2 - 0.3*x + 0.2,
                    -x*x/8 - 0.075*x + 0.2)


def exact_deriv(x):
    x = np.asarray(x, dtype=float)
    return np.where(x < 0, -x - 0.3, -x/4 - 0.075)


def solve_tridiagonal(subdiag, diag, rhs):
    """Algorytm Thomasa: macierz symetryczna trojdiagonalna."""
    b, d, c = diag.copy(), rhs.copy(), subdiag.copy()
    for i in range(1, len(b)):
        fac = c[i - 1] / b[i - 1]
        b[i] -= fac * c[i - 1]
        d[i] -= fac * d[i - 1]
    x = np.empty_like(d)
    x[-1] = d[-1] / b[-1]
    for i in range(len(x)-2, -1, -1):
        x[i] = (d[i] - c[i] * x[i+1]) / b[i]
    return x


def fem_solution(n_per_half):
    """Liniowe FEM dla -(a u')'=1, u(-1)=u(1)=0.

    Punkt 0 jest wezelm siatki. a=1 po lewej i a=4 po prawej.
    """
    nodes = np.linspace(-1., 1., 2 * n_per_half + 1)
    h = nodes[1] - nodes[0]
    a_elem = np.where((nodes[:-1] + nodes[1:]) / 2 < 0, 1., 4.)
    diag = (a_elem[:-1] + a_elem[1:]) / h
    subdiag = -a_elem[1:-1] / h
    rhs = np.full(len(nodes)-2, h)
    sol = np.r_[0., solve_tridiagonal(subdiag, diag, rhs), 0.]
    assert np.max(np.abs(sol - exact(nodes))) < 5e-10
    return nodes, sol, a_elem


def errors(nodes, sol):
    z, w = np.polynomial.legendre.leggauss(5)
    l2sq = h1sq = 0.
    for j in range(len(nodes) - 1):
        x0, x1 = nodes[j:j+2]
        h = x1 - x0
        xx = (x0+x1)/2 + h*z/2
        approx = sol[j] + (sol[j+1]-sol[j])*(xx-x0)/h
        slope = (sol[j+1]-sol[j])/h
        l2sq += (h/2) * np.dot(w, (approx-exact(xx))**2)
        h1sq += (h/2) * np.dot(w, (slope-exact_deriv(xx))**2)
    return np.sqrt(l2sq), np.sqrt(h1sq)


def save(fig, name):
    fig.update_layout(template="plotly_white", font=dict(size=13),
                      hovermode="x unified")
    fig.write_html(OUT / name, include_plotlyjs=True, full_html=True)


def main():
    hs, l2errors, h1errors = [], [], []
    fig_sol = go.Figure()
    x = np.linspace(-1, 1, 1201)
    fig_sol.add_trace(go.Scatter(x=x, y=exact(x), mode="lines",
                                 name="Rozwiazanie dokladne"))
    fig_flux = go.Figure()
    fig_flux.add_trace(go.Scatter(x=x, y=x + 0.3, mode="lines",
                                  name="Dokladny strumien q=-a u'"))
    for N in [4, 8, 16, 32, 64, 128]:
        nodes, vals, a_elem = fem_solution(N)
        h = nodes[1]-nodes[0]
        l2, h1 = errors(nodes, vals)
        assert np.isclose(l2, h*h * np.sqrt(17/1920), rtol=5e-8, atol=5e-11)
        assert np.isclose(h1, h * np.sqrt(17/192), rtol=5e-8, atol=5e-11)
        hs.append(h)
        l2errors.append(l2)
        h1errors.append(h1)
        if N in (4, 16):
            fig_sol.add_trace(go.Scatter(x=nodes, y=vals, mode="lines+markers",
                                         name=f"FEM, N={N} na polowe"))
        if N in (4, 16):
            mid = (nodes[:-1]+nodes[1:])/2
            qfem = -a_elem*np.diff(vals)/h
            fig_flux.add_trace(go.Scatter(x=mid, y=qfem, mode="markers",
                                          name=f"q_h w srodkach, N={N}"))
        print(f"N={N:3d} h={h:.5f} L2={l2:.9e} H1-semi={h1:.9e}")
    fig_sol.add_vline(x=0, line_dash="dash", annotation_text="Interfejs")
    fig_sol.update_layout(title="Transmisja przez dwa osrodki: rozwiazanie slabe",
                          xaxis_title="x", yaxis_title="u(x)")
    save(fig_sol, "pde13_transmisja_rozwiazanie.html")
    fig_flux.add_vline(x=0, line_dash="dash", annotation_text="Interfejs")
    fig_flux.update_layout(title="Strumien ciagly w rozwiazaniu dokladnym",
                           xaxis_title="x", yaxis_title="q=-a u'(x)")
    save(fig_flux, "pde13_transmisja_strumien.html")
    fig_error = go.Figure()
    fig_error.add_trace(go.Scatter(x=hs, y=l2errors, mode="lines+markers",
                                   name="Blad L2, rzad 2"))
    fig_error.add_trace(go.Scatter(x=hs, y=h1errors, mode="lines+markers",
                                   name="Blad seminormy H1, rzad 1"))
    fig_error.update_layout(title="Zbieznosc FEM dla rownania eliptycznego",
                            xaxis_title="h", yaxis_title="Blad",
                            xaxis_type="log", yaxis_type="log")
    save(fig_error, "pde13_transmisja_zbieznosc.html")
    p2 = np.polyfit(np.log(hs), np.log(l2errors), 1)[0]
    p1 = np.polyfit(np.log(hs), np.log(h1errors), 1)[0]
    assert np.isclose(p2, 2., atol=1e-6) and np.isclose(p1, 1., atol=1e-6)
    print(f"Rzedy: L2={p2:.6f}, H1-semi={p1:.6f}; testy zaliczone.")


if __name__ == "__main__":
    main()
\end{lstlisting}

\section{Zasady maksimum, nierówność Harnacka i lemat Hopfa — Plotly}
Kod oblicza harmoniczne przedłużenie danych brzegowych w dysku,
weryfikuje nierówność Harnacka oraz pochodną normalną w przykładzie
ilustrującym lemat Hopfa. Porównanie jądra Poissona z rozwiązaniem
analitycznym jest jednocześnie testem poprawności normalizacji jądra.
\begin{lstlisting}[style=pythonpde,caption={Jądro Poissona, nierówność Harnacka i lemat Hopfa (Plotly)},label={code:pde14-plotly}]
"""PDE, rozdzial 14: jadro Poissona, Harnack i Hopf.

Instalacja: pip install numpy plotly
Uruchomienie: python demo_maksimum_harnack_plotly.py
Wykresy HTML zawieraja Plotly offline.
"""
from pathlib import Path
import numpy as np
import plotly.graph_objects as go

OUT = Path(__file__).resolve().parent
TWO_PI = 2.0 * np.pi


def save(fig, filename):
    fig.update_layout(template="plotly_white", font=dict(size=13))
    fig.write_html(OUT / filename, include_plotlyjs=True, full_html=True)


def boundary(phi):
    return 0.5 + np.cos(2.0 * phi) + 0.3 * np.sin(3.0 * phi)


def poisson_disk(r, theta, phi):
    # Jadrowa gestosc wzgledem miary dlugosci na okregu jednostkowym.
    return (1.0 - r * r) / (TWO_PI * (1.0 + r * r - 2.0 * r * np.cos(theta - phi)))


def harmonic_exact(r, theta):
    return 0.5 + r**2 * np.cos(2.0 * theta) + 0.3 * r**3 * np.sin(3.0 * theta)


def main():
    phi = np.linspace(0.0, TWO_PI, 2048, endpoint=False)
    theta = np.linspace(0.0, TWO_PI, 160)
    g = boundary(phi)
    fig = go.Figure()
    max_error = 0.0
    for r in [0.0, 0.3, 0.6, 0.85]:
        p = poisson_disk(r, theta[:, None], phi[None, :])
        vals = TWO_PI * np.mean(p * g[None, :], axis=1)
        exact = harmonic_exact(r, theta)
        error = np.max(np.abs(vals - exact))
        max_error = max(max_error, error)
        fig.add_trace(go.Scatter(x=theta, y=vals, mode="lines", name=f"r={r:g}"))
        assert error < 1e-9, (r, error)
    fig.update_layout(title="Przedluzenie harmoniczne przez jadro Poissona",
                      xaxis_title="kat theta", yaxis_title="u(r, theta)")
    save(fig, "pde14_poisson_plotly.html")
    print(f"Poisson: maksymalny blad dyskretyzacji = {max_error:.3e}")

    # Jawna funkcja harmoniczna: u(x,y)=1+0.6*x w B_1(0).
    r = 0.5
    t = np.linspace(0.0, TWO_PI, 500)
    vals = 1.0 + 0.6 * r * np.cos(t)
    ratio = np.max(vals) / np.min(vals)
    lower = (1.0 - r**2) / (1.0 + r)**2
    upper = 1.0 / (1.0 - r)**2
    assert np.min(vals) >= lower and np.max(vals) <= upper
    fig = go.Figure()
    fig.add_trace(go.Scatter(x=t, y=vals, name="u(0.5 cos(theta), 0.5 sin(theta))"))
    fig.add_hline(y=lower, line_dash="dash", annotation_text="dolne Harnack")
    fig.add_hline(y=upper, line_dash="dash", annotation_text="gorne Harnack")
    fig.update_layout(title=f"Harnack: rzeczywisty sup/inf = {ratio:.4f}",
                      xaxis_title="kat theta", yaxis_title="u")
    save(fig, "pde14_harnack_plotly.html")
    print(f"Harnack: sup/inf = {ratio:.6f}, zakres teoretyczny = [{lower:.3f}, {upper:.3f}]")

    # Dodatnia funkcja harmoniczna u(x,y)=1-x; w x0=(1,0) minimum 0.
    x = np.linspace(0.75, 1.0, 150)
    u = 1.0 - x
    fig = go.Figure()
    fig.add_trace(go.Scatter(x=x, y=u, name="u(x,0)=1-x"))
    fig.add_trace(go.Scatter(x=x, y=(1.0-x), mode="markers", name="styczna w x0", marker=dict(size=3)))
    fig.update_layout(title="Lemat Hopfa: u(1,0)=0, pochodna normalna = -1",
                      xaxis_title="x (promien do brzegu)", yaxis_title="u(x,0)")
    save(fig, "pde14_hopf_plotly.html")
    assert np.allclose(np.gradient(u, x), -1.0, atol=1e-10)
    print("Hopf: pochodna zewnetrzna = -1 (blad < 1e-10)")


if __name__ == "__main__":
    main()
\end{lstlisting}

\section{Ilorazy różnicowe i skok pochodnej — Plotly}
Przykład pokazuje, jak skok pochodnej prowadzi do nieograniczonej normy dyskretnych drugich pochodnych.
\begin{lstlisting}[style=pythonpde,caption={Brak regularnosci H2 dla transmisji (Plotly)},label={code:pde15-plotly}]
"""PDE, rozdzial 15: H2-regularnosc i transmisja (Plotly).
Instalacja: pip install numpy plotly
"""
from pathlib import Path
import numpy as np
import plotly.graph_objects as go

OUT = Path(__file__).resolve().parent

def save(fig, name):
    fig.update_layout(template="plotly_white", font=dict(size=13))
    fig.write_html(OUT / name, include_plotlyjs=True, full_html=True)

def transmission(x):
    return np.where(x <= 0.5, x, 0.5 + 0.5*(x-0.5))

def main():
    x = np.linspace(0., 1., 1001)
    u = transmission(x)
    fig = go.Figure()
    fig.add_trace(go.Scatter(x=x, y=u, name="u: ciagla, z zalamaniem"))
    fig.add_trace(go.Scatter(x=x, y=x*(1-x), name="v: gladka"))
    fig.add_vline(x=0.5, line_dash="dot", annotation_text="granica osrodkow")
    fig.update_layout(title="Regularnosc rozwiazan: gladkie i transmisyjne",
                      xaxis_title="x", yaxis_title="wartosc")
    save(fig, "pde15_regularnosc_plotly.html")

    Ns = np.array([20, 40, 80, 160, 320, 640])
    hs = 1.0/Ns
    norms = []
    for N in Ns:
        xgrid = np.linspace(0, 1, N+1)
        uh = transmission(xgrid)
        second = (uh[2:] - 2*uh[1:-1] + uh[:-2]) / (1/N)**2
        norms.append(np.sqrt(np.sum(second**2) / N))
    norms = np.array(norms)
    reference = 0.5 / np.sqrt(hs)
    assert np.allclose(norms, reference, rtol=1e-10, atol=1e-10)
    slope = np.polyfit(np.log(hs), np.log(norms), 1)[0]
    assert abs(slope + .5) < 1e-10
    fig = go.Figure()
    fig.add_trace(go.Scatter(x=hs, y=norms, mode="lines+markers", name="dyskretna norma H2"))
    fig.add_trace(go.Scatter(x=hs, y=reference, mode="lines", name="0.5 / sqrt(h)"))
    fig.update_layout(title=f"Brak H2: norma drugiej roznicy ~ h^({slope:.2f})",
                      xaxis_type="log", yaxis_type="log",
                      xaxis_title="h", yaxis_title="norma L2 drugiej roznicy")
    save(fig, "pde15_skok_h2_plotly.html")
    print(f"Nachylenie log-log: {slope:.8f} (oczekiwane -0.5)")

if __name__ == "__main__":
    main()
\end{lstlisting}

\section[Drugie pochodne i ostrość oszacowań Schaudera — Plotly]{Drugie pochodne, brak oszacowania w $L^\infty$ i ostrość oszacowań Schaudera — Plotly}
Poniższy eksperyment na torusie dwuwymiarowym ilustruje kontrast
między normą średniokwadratową a normą supremum Hessianu.
Jest to eksperyment periodyczny, nie dowód twierdzenia dla $\R^d$.
Na końcu kod przedstawia dokładny przykład ostrości wykładnika Höldera.
\begin{lstlisting}[style=pythonpde,caption={Rachunek spektralny dla Hessianu i wykresy ilustrujące oszacowania Schaudera (Plotly)},label={code:pde16-plotly}]
"""Rozdzial 16: regularnosc W2p i oszacowania Schaudera; interaktywne wykresy Plotly.
Instalacja: pip install numpy plotly
Uruchomienie: python demo_calderon_schauder_plotly.py
"""
from pathlib import Path
import numpy as np
import plotly.graph_objects as go

OUT = Path(__file__).resolve().parent


def save(fig, name):
    fig.update_layout(template='plotly_white', font=dict(size=13))
    fig.write_html(OUT / name, include_plotlyjs=True, full_html=True)


def mixed_hessian(N):
    # Kwadrat w torusie [0,2pi)^2; z jego funkcji wskaznikowej
    # usuwamy srednia, by rozwiazanie -Delta u=f istnialo okresowo.
    x = 2 * np.pi * np.arange(N) / N
    square = ((x[:, None] < np.pi) & (x[None, :] < np.pi)).astype(float)
    f = square - square.mean()
    k = np.fft.fftfreq(N, d=1.0/N)  # calkowite czestosci
    K1, K2 = np.meshgrid(k, k, indexing='ij')
    k2 = K1*K1 + K2*K2
    m = np.zeros_like(k2, dtype=float)
    np.divide(-K1*K2, k2, out=m, where=k2 != 0)
    d12u = np.fft.ifft2(np.fft.fft2(f)*m).real
    return x, f, d12u


def demo_fft():
    N = 256
    x, f, d12u = mixed_hessian(N)
    fig = go.Figure(data=go.Heatmap(x=x, y=x, z=d12u.T,
                                     colorbar=dict(title='u_xy')))
    fig.update_layout(title='Drugie pochodne rozwiazania -Delta u=f, dane naroznikowe',
                      xaxis_title='x', yaxis_title='y', yaxis_scaleanchor='x')
    save(fig, 'pde16_fft_hessjan.html')

    Ns = [32, 64, 128, 256, 512]
    peaks = []
    rms = []
    for n in Ns:
        _, fn, hn = mixed_hessian(n)
        peaks.append(float(np.max(np.abs(hn))))
        rms.append(float(np.mean(hn**2)**0.5))
        # Mnoznik Hessianu ma modul <= 1 w normie L2:
        assert rms[-1] <= float(np.mean(fn**2)**0.5) + 1e-11
    fig = go.Figure()
    fig.add_trace(go.Scatter(x=Ns, y=peaks, mode='lines+markers', name='maksimum |u_xy|'))
    fig.add_trace(go.Scatter(x=Ns, y=rms, mode='lines+markers', name='norma sredniokwadratowa'))
    fig.update_layout(title='Wzrost maksimum przy zachowanym oszacowaniu L2',
                      xaxis_title='liczba wezlow N', xaxis_type='log',
                      yaxis_title='wartosc normy')
    save(fig, 'pde16_fft_normy.html')
    print('N / maksimum / RMS')
    for n, peak, r in zip(Ns, peaks, rms):
        print(f'{n:4d} / {peak:.5f} / {r:.5f}')
    assert peaks[-1] > peaks[0], 'Oczekiwano wzrostu maksimum przy narozniku'


def demo_holder():
    x = np.linspace(-1, 1, 1501)
    alpha = 0.6
    u = -np.abs(x)**(alpha+2)/((alpha+1)*(alpha+2))
    f = np.abs(x)**alpha
    fig = go.Figure()
    fig.add_trace(go.Scatter(x=x, y=u, name='u: rozwiazanie -u\'\'=f'))
    fig.add_trace(go.Scatter(x=x, y=f, name='f=|x|^alpha = -u\'\''))
    fig.update_layout(title=f'Ostrosc oszacowan Schaudera, alpha={alpha}',
                      xaxis_title='x', yaxis_title='wartosc')
    save(fig, 'pde16_holder_ostrosc.html')
    h = np.logspace(-4, -1, 30)
    beta = 0.8
    qa = h**alpha / h**alpha
    qb = h**alpha / h**beta
    fig = go.Figure()
    fig.add_trace(go.Scatter(x=h, y=qa, name='iloraz Holdera dla beta=alpha'))
    fig.add_trace(go.Scatter(x=h, y=qb, name='iloraz Holdera dla beta>alpha'))
    fig.update_layout(title='Wyzszy wykladnik Holdera jest nieosiagalny',
                      xaxis_type='log', yaxis_type='log',
                      xaxis_title='odleglosc h', yaxis_title='iloraz')
    save(fig, 'pde16_holder_ilorazy.html')
    assert np.allclose(qa, 1)
    assert qb[0] > qb[-1]


if __name__ == '__main__':
    demo_fft()
    demo_holder()
\end{lstlisting}

\section{Regularność De Giorgiego--Nasha--Mosera: przykład radialny — Plotly}
\label{sec:pde17-code}
Program wizualizuje słabe rozwiązanie o optymalnej regularności
$C^{0,\gamma}$, macierz współczynników oraz ilorazy Höldera.
Dodatkowo sprawdza wartości maksimum i minimum przykładowej
funkcji harmonicznej. Jest to eksperyment ilustracyjny,
nie dowód twierdzenia o regularności przy mierzalnych współczynnikach.
\begin{lstlisting}[style=pythonpde,caption={Mierzalne współczynniki, regularność Höldera i Harnack (Plotly)},label={code:pde17-plotly}]
"""PDE rozdzial 17: ostrosc Holdera przy mierzalnych wspolczynnikach.

Zaleznosci: pip install numpy plotly
Uruchomienie: python demo_degiorgi_moser_plotly.py
Wyniki: pde17_profil_holder.html, pde17_macierz.html,
        pde17_ilorazy_holdera.html, pde17_harnack_affine.html
"""
from pathlib import Path
import numpy as np
import plotly.graph_objects as go

OUT = Path(__file__).resolve().parent


def store(fig: go.Figure, filename: str) -> None:
    fig.update_layout(template="plotly_white", font=dict(size=13))
    fig.write_html(OUT / filename, include_plotlyjs=True, full_html=True)


def radial_example(gamma: float, n: int = 301):
    xx = np.linspace(-1.0, 1.0, n)
    x, y = np.meshgrid(xx, xx, indexing="xy")
    radius = np.hypot(x, y)
    theta_cos = np.divide(x, radius, out=np.zeros_like(x), where=radius > 0)
    u = radius**gamma * theta_cos
    u[radius > 1] = np.nan
    # A = gamma^2 I + (1-gamma^2) e_r \otimes e_r
    e1 = np.divide(x, radius, out=np.zeros_like(x), where=radius > 0)
    e2 = np.divide(y, radius, out=np.zeros_like(y), where=radius > 0)
    a11 = gamma**2 + (1.0-gamma**2) * e1**2
    a22 = gamma**2 + (1.0-gamma**2) * e2**2
    a12 = (1.0-gamma**2) * e1 * e2
    return xx, x, y, radius, u, a11, a12, a22


def main():
    gamma = 0.55
    xx, x, y, radius, u, a11, a12, a22 = radial_example(gamma)
    fig = go.Figure(go.Heatmap(x=xx, y=xx, z=u, colorscale="RdBu",
                               zmid=0, colorbar=dict(title="u")))
    fig.update_layout(title=f"Rozwiazanie r^gamma cos(theta), gamma={gamma}",
                      xaxis_title="x", yaxis_title="y", yaxis_scaleanchor="x")
    store(fig, "pde17_profil_holder.html")

    fig = go.Figure()
    fig.add_trace(go.Heatmap(x=xx, y=xx, z=np.where(radius <= 1, a12, np.nan),
                             colorscale="RdBu", zmid=0, colorbar=dict(title="a12")))
    fig.update_layout(title="Mierzalna macierz A: wspolczynnik a12",
                      xaxis_title="x", yaxis_title="y", yaxis_scaleanchor="x")
    store(fig, "pde17_macierz.html")

    r = np.logspace(-4, -1, 90)
    fig = go.Figure()
    for beta in [0.35, gamma, 0.75]:
        quotient = r**(gamma-beta)
        fig.add_trace(go.Scatter(x=r, y=quotient, mode="lines",
                                 name=f"beta={beta}"))
    fig.update_layout(title="Ilorazy Holdera w zerze: r^(gamma-beta)",
                      xaxis_title="r", yaxis_title="iloraz", xaxis_type="log",
                      yaxis_type="log")
    store(fig, "pde17_ilorazy_holdera.html")
    assert np.all(np.diff(r**(gamma-0.75)) < 0)

    rr = np.linspace(0.0, 0.95, 100)
    # harmonic positive v(x,y)=2+0.4x on B_1
    sup = 2.0 + 0.4 * rr
    inf = 2.0 - 0.4 * rr
    ratio = sup/inf
    assert np.max(ratio) < 1.5
    fig = go.Figure()
    fig.add_trace(go.Scatter(x=rr, y=sup, name="sup na B_r"))
    fig.add_trace(go.Scatter(x=rr, y=inf, name="inf na B_r"))
    fig.add_trace(go.Scatter(x=rr, y=ratio, name="sup/inf"))
    fig.update_layout(title="Przykladowa dodatnia funkcja harmoniczna i Harnack",
                      xaxis_title="promien r", yaxis_title="wartosc")
    store(fig, "pde17_harnack_affine.html")
    print("Zweryfikowano cztery wykresy i prosty test Harnacka.")


if __name__ == "__main__":
    main()
\end{lstlisting}


\section{Widmo eliptyczne: wartości własne, Weyl i ślad cieplny — Plotly}
\label{sec:pde18-code}
Poniższy program oblicza dyskretne wartości własne jednowymiarowego
laplasjanu Dirichleta, porównuje funkcję zliczającą kwadratowego
laplasjanu różnicowego z wyrazem głównym prawa Weyla oraz
wizualizuje ślad półgrupy cieplnej. Ostatnie porównanie wykorzystuje
jawne wartości własne ciągłego laplasjanu na kwadracie.
Przy ustalonej siatce nie należy interpretować końca widma
macierzy jako asymptotyki wysokich częstotliwości operatora ciągłego.
\begin{lstlisting}[style=pythonpde,caption={Wartości własne, prawo Weyla i ślad cieplny (Plotly)},label={code:pde18-plotly}]
"""PDE, rozdzial 18: min-max, Weyl, slad cieplny.

Zaleznosci: pip install numpy scipy plotly
Uruchom: python demo_widmo_weyl_plotly.py
Wyniki: pde18_wartosci_wlasne.html, pde18_funkcja_zliczajaca.html,
        pde18_slad_cieplny.html
"""
from pathlib import Path
import numpy as np
import plotly.graph_objects as go
from scipy.linalg import eigh_tridiagonal

OUT = Path(__file__).resolve().parent


def save(fig, filename):
    fig.update_layout(template='plotly_white', font=dict(size=13),
                      margin=dict(l=60, r=35, t=65, b=65))
    fig.write_html(OUT / filename, include_plotlyjs=True, full_html=True)


def dirichlet_interval():
    ns = np.array([12, 20, 32, 50, 80, 128, 200])
    errors = []
    for n in ns:
        h = 1.0 / (n + 1)
        vals = eigh_tridiagonal(np.full(n, 2 / h**2),
                                np.full(n - 1, -1 / h**2),
                                select='i', select_range=(0, 0),
                                eigvals_only=True)
        errors.append(abs(vals[0] - np.pi**2))
    err = np.asarray(errors)
    h = 1 / (ns + 1)
    fig = go.Figure()
    fig.add_trace(go.Scatter(x=h, y=err, mode='lines+markers',
                             name='Blad pierwszej wartosci wlasnej'))
    fig.add_trace(go.Scatter(x=h, y=err[-1] * (h/h[-1])**2,
                             mode='lines', name='Linia odniesienia: h^2'))
    fig.update_layout(title='Zbieznosc wartosci wlasnej laplasjanu (Dirichlet)',
                      xaxis_type='log', yaxis_type='log',
                      xaxis_title='Krok h', yaxis_title='Blad bezwzgledny')
    save(fig, 'pde18_wartosci_wlasne.html')
    slope = np.polyfit(np.log(h), np.log(err), 1)[0]
    assert 1.95 < slope < 2.05, f'Oczekiwany rzad 2, otrzymano {slope}'
    print('Rzad zbieznosci wartosci wlasnej:', round(slope, 6))


def weyl_square():
    m = 100
    h = 1 / (m + 1)
    n = np.arange(1, m+1)
    one_d = 4 * np.sin(np.pi * n/(2*(m+1)))**2 / h**2
    two_d = np.sort((one_d[:, None] + one_d[None, :]).ravel())
    cutoffs = np.linspace(0, 2400, 450)
    counts = np.searchsorted(two_d, cutoffs, side='right')
    leading = cutoffs / (4*np.pi)  # |Omega|=1, d=2
    fig = go.Figure()
    fig.add_trace(go.Scatter(x=cutoffs, y=counts, mode='lines',
                             line_shape='hv', name='Dyskretne zliczanie modow'))
    fig.add_trace(go.Scatter(x=cutoffs, y=leading, mode='lines',
                             name='Wyraz glowny Weyla: Lambda/(4 pi)'))
    fig.update_layout(title='Funkcja zliczajaca na kwadracie jednostkowym',
                      xaxis_title='Prog Lambda', yaxis_title='Liczba modow')
    save(fig, 'pde18_funkcja_zliczajaca.html')
    # m=100 makes low modes accurate; at high frequencies the grid
    # is fixed and hence Weyl is only illustrative.
    print('N_dyskretne(1000)=', np.searchsorted(two_d, 1000, side='right'))


def heat_trace():
    j = np.arange(1, 180)
    eigs = np.pi**2 * (j[:, None]**2 + j[None, :]**2)
    t = np.logspace(-3.3, -1.2, 170)
    exact = np.array([np.exp(-tau*eigs).sum() for tau in t])
    leading = 1 / (4*np.pi*t)  # |Omega|=1, d=2
    fig = go.Figure()
    fig.add_trace(go.Scatter(x=t, y=exact, mode='lines',
                             name='Slad cieplny (suma modow)'))
    fig.add_trace(go.Scatter(x=t, y=leading, mode='lines',
                             name='Asymptotyka 1/(4 pi t)'))
    fig.update_layout(title='Slad cieplny laplasjanu Dirichleta na kwadracie',
                      xaxis_type='log', yaxis_type='log',
                      xaxis_title='Czas t', yaxis_title='Tr exp(t Delta_D)')
    save(fig, 'pde18_slad_cieplny.html')
    ratio = exact[0] / leading[0]
    assert 0.8 < ratio < 1.0
    print('Stosunek sladu do wyrazu glownego dla malego t:', round(ratio, 6))


if __name__ == '__main__':
    dirichlet_interval()
    weyl_square()
    heat_trace()
\end{lstlisting}



\section{Pojemność kondensatora kulistego — Plotly}
\label{sec:pde19-code}
Poniższy program porównuje pojemności względne
$\operatorname{cap}_{2,B_1}(\overline B_r)$ w wymiarach $2$ i $3$.
Każda krzywa jest obliczana z dokładnego wzoru
\eqref{eq:pde19-ball-exact}; program nie zastępuje dowodu
minimalności energii ani nie numeryzuje równania Laplace'a.
Wykresy zapisywane są jako interaktywne pliki HTML.
\begin{lstlisting}[style=pythonpde,caption={Pojemność Sobolewa i tempo zaniku pojemności punktu (Plotly)},label={code:pde19-plotly}]
"""PDE, rozdzial 19: pojemnosc kondensatora kulistego.
Wymaga: numpy, plotly. Wyniki: interaktywne pliki HTML.
"""
from pathlib import Path
import numpy as np
import plotly.graph_objects as go

OUT = Path(__file__).resolve().parent
R = 1.0
r = np.geomspace(1e-6, 0.9, 240)
cap2 = 2 * np.pi / np.log(R / r)
cap3 = 4 * np.pi / (1 / r - 1 / R)


def save(fig, name):
    fig.update_layout(template="plotly_white", font=dict(size=13))
    fig.write_html(OUT / name, include_plotlyjs=True, full_html=True)


def condenser_plot(cap, d, name):
    fig = go.Figure()
    fig.add_trace(go.Scatter(x=r, y=cap, mode="lines",
                             name=f"pojemnosc d={d}"))
    fig.update_layout(title=f"Pojemnosc wzgledna kuli w B_1, d={d}",
                      xaxis_title="promien r", yaxis_title="energia minimalna")
    save(fig, name)


def comparison():
    fig = go.Figure()
    fig.add_trace(go.Scatter(x=r, y=cap2, name="d=2 (logarytm)"))
    fig.add_trace(go.Scatter(x=r, y=cap3, name="d=3 (potega)"))
    fig.update_layout(title="Zanik pojemnosci punktu, rozne tempa",
                      xaxis_type="log", yaxis_type="log",
                      xaxis_title="promien r", yaxis_title="pojemnosc")
    save(fig, "pde19_pojemnosc_porownanie.html")


def test_exact_values():
    assert np.all(np.diff(cap2) > 0)
    assert np.all(np.diff(cap3) > 0)
    assert np.isclose(2 * np.pi / np.log(2), 9.064720283654388)
    assert np.isclose(4 * np.pi / (2 - 1), 4 * np.pi)
    print("d=2, r=1/2:", 2 * np.pi / np.log(2))
    print("d=3, r=1/2:", 4 * np.pi)
    print("d=2, r=1e-6:", cap2[0])
    print("d=3, r=1e-6:", cap3[0])


if __name__ == "__main__":
    test_exact_values()
    condenser_plot(cap2, 2, "pde19_pojemnosc_2d.html")
    condenser_plot(cap3, 3, "pde19_pojemnosc_3d.html")
    comparison()
\end{lstlisting}


\section{Paraboliczne równania ewolucyjne: wsteczny Euler i stabilność — Plotly}
\label{sec:pde20-code}
Eksperyment porównuje niejawny schemat różnicowy z rozwiązaniem dokładnym
na odcinku, testuje rząd zbieżności przy $\tau\asymp h^2$ oraz
ilustruje tłumienie poszczególnych modów. W algorytmie stosujemy
dyskretną transformatę sinusową, która diagonalizuje laplasjan
z jednorodnymi warunkami Dirichleta. Jest to implementacja
wstecznego Eulera, a nie schematu dokładnej ewolucji macierzy.
\begin{lstlisting}[style=pythonpde,caption={Niejawny schemat cieplny, zbieżność i tłumienie modów (Plotly)},label={code:pde20-plotly}]
"""Iteracja 20, rownanie ciepa na (0,1) z Dirichletem.

Wsteczny Euler i centralny laplasjan roznicowy,
diagonalizacja macierzy przez dyskretna transformate sinusowa.
Instalacja: python -m pip install numpy scipy plotly
Uruchomienie: python demo_paraboliczne_plotly.py
"""
from pathlib import Path
import numpy as np
import plotly.graph_objects as go
from scipy.fft import dst, idst

OUT = Path(__file__).resolve().parent
T = 0.035


def initial(x):
    return np.sin(np.pi * x) + 0.5 * np.sin(3 * np.pi * x)


def exact(x, t):
    return np.exp(-np.pi**2 * t) * np.sin(np.pi*x) + 0.5 * np.exp(-9*np.pi**2*t) * np.sin(3*np.pi*x)


def backward_euler(n, at=T, cfl=0.32):
    """Rozwiazanie po nstep krokach: kazdy krok jest niejawny."""
    h = 1 / (n + 1)
    nstep = max(1, int(np.ceil(at / (cfl*h*h))))
    tau = at / nstep
    x = np.arange(1, n+1) * h
    j = np.arange(1, n+1)
    eig = 4 / h**2 * np.sin(j*np.pi/(2*(n+1)))**2
    coeffs = dst(initial(x), type=1, norm='ortho')
    ampl = (1 + tau*eig)**(-nstep)
    uh = idst(coeffs*ampl, type=1, norm='ortho')
    return x, uh, tau, eig


def save(fig, file):
    fig.update_layout(template='plotly_white', font=dict(size=13),
                      margin=dict(l=70, r=35, b=60, t=75))
    fig.write_html(OUT/file, include_plotlyjs=True, full_html=True)


def profiles():
    n=80
    h = 1/(n+1)
    x = np.arange(1,n+1)*h
    fig = go.Figure()
    for t in [0.0, 0.004, 0.012, T]:
        if t == 0:
            u=initial(x)
        else:
            _, u, _, _ = backward_euler(n, at=t)
        fig.add_trace(go.Scatter(x=x, y=u, mode='lines', name=f't={t:g}'))
    fig.update_layout(title='Rownanie ciepa: wygadzanie i zanikanie modow',
                      xaxis_title='x',yaxis_title='u(t,x)')
    save(fig, 'pde20_profile_plotly.html')
    # Zanik nie zwieksza normy w l2 dla danych bez wymuszenia.
    assert np.linalg.norm(backward_euler(n)[1]) <= np.linalg.norm(initial(x))


def convergence():
    sizes=np.array([12,18,27,40,60,90,135])
    errs=[]
    for n in sizes:
        x, u, _, _ = backward_euler(int(n))
        h = 1/(n+1)
        errs.append(np.sqrt(h*np.sum((u-exact(x,T))**2)))
    hs=1/(sizes+1)
    errs=np.array(errs)
    slope=np.polyfit(np.log(hs),np.log(errs),1)[0]
    print('Wsteczny Euler / laplasjan roznicowy: rzad bedu',round(slope,5))
    assert 1.85 < slope < 2.12, f'Nieoczekiwany rzad: {slope}'
    fig=go.Figure()
    fig.add_trace(go.Scatter(x=hs,y=errs,mode='lines+markers',name='Bad w L2_h'))
    fig.add_trace(go.Scatter(x=hs,y=errs[-1]*(hs/hs[-1])**2,name='Odniesienie: h2'))
    fig.update_layout(title=f'Zbieznosc dla 0.32h2, rzad {slope:.3f}',
                      xaxis_type='log',yaxis_type='log',
                      xaxis_title='h',yaxis_title='||U-u||_(L2_h)')
    save(fig, 'pde20_zbieznosc_plotly.html')


def damping():
    z=np.logspace(-3,2,500)
    fig=go.Figure()
    fig.add_trace(go.Scatter(x=z,y=np.exp(-z),name='Dokadny: exp(-z)'))
    fig.add_trace(go.Scatter(x=z,y=1/(1+z),name='Wsteczny Euler: 1/(1+z)'))
    fig.add_trace(go.Scatter(x=z,y=np.abs(1-z),name='Jawny Euler: |1-z|'))
    fig.update_layout(title='Tumienie jednego modu; z=',
                      xaxis_title='z',yaxis_title='wspoczynnik amplitudy',
                      xaxis_type='log',yaxis_type='log')
    save(fig,'pde20_tlumienie_plotly.html')
    assert np.all(1/(1+z)<=1)
    assert np.any(np.abs(1-z)>1)


if __name__ == '__main__':
    profiles()
    convergence()
    damping()
\end{lstlisting}


\section{Układ symetryczno-hiperboliczny: schemat pod prąd — Plotly}
\label{sec:pde21-code}
Program oblicza ewolucję dwóch modów o różnych prędkościach,
porównuje rozwiązanie numeryczne z dokładnym,
sprawdza nierosnącą energię dyskretną i wyznacza rząd
zbieżności jawnego schematu pod prąd. Rozdzielenie macierzy
$A=A_++A_-$ odbywa się przez diagonalizację ortogonalną,
a warunek CFL ma postać $\tau\max_k|a_k|/h\le1$.
\begin{lstlisting}[style=pythonpde,caption={Hiperboliczny układ, energia i zbieżność: Plotly},label={code:pde21-plotly}]
"""Rozdzial 21: dwa przemieszczajace sie mody ukadu symetrycznego.
Zaleznosci: numpy, plotly; uruchomienie: python demo_hiperboliczne_plotly.py
Schemat pod prad z rozdzieleniem macierzy A=A_++A_-.
"""
from pathlib import Path
import numpy as np
import plotly.graph_objects as go

OUT = Path(__file__).resolve().parent
THETA = 0.48
Q = np.array([[np.cos(THETA), -np.sin(THETA)],
              [np.sin(THETA), np.cos(THETA)]])
SPEEDS = np.array([-1.0, 0.6])
A = Q @ np.diag(SPEEDS) @ Q.T
A_PLUS = Q @ np.diag(np.maximum(SPEEDS, 0.0)) @ Q.T
A_MINUS = Q @ np.diag(np.minimum(SPEEDS, 0.0)) @ Q.T


def wrapped_gauss(x, center, width):
    y = (x - center + 0.5) % 1.0 - 0.5
    return np.exp(-(y / width) ** 2)


def initial(x):
    w = np.vstack([wrapped_gauss(x, 0.25, 0.065),
                   -0.8 * wrapped_gauss(x, 0.71, 0.075)])
    return Q @ w


def exact(x, t):
    w = np.vstack([wrapped_gauss(x - SPEEDS[0] * t, 0.25, 0.065),
                   -0.8 * wrapped_gauss(x - SPEEDS[1] * t, 0.71, 0.075)])
    return Q @ w


def upwind(n, T=0.35, cfl=0.8, record=False):
    h = 1.0 / n
    x = np.arange(n) * h
    steps = int(np.ceil(T * np.max(abs(SPEEDS)) / (cfl * h)))
    dt = T / steps
    U = initial(x)
    en = [h * np.sum(U * U)]
    for _ in range(steps):
        left = U - np.roll(U, 1, axis=1)
        right = np.roll(U, -1, axis=1) - U
        U = U - dt / h * (A_PLUS @ left + A_MINUS @ right)
        if record:
            en.append(h * np.sum(U * U))
    ref = exact(x, T)
    err = np.sqrt(h * np.sum((U - ref) ** 2))
    return x, U, ref, err, np.array(en), dt, h


def export(fig, name):
    fig.update_layout(template='plotly_white', font=dict(size=13),
                      margin=dict(l=75, r=35, b=55, t=75))
    fig.write_html(OUT / name, include_plotlyjs=True, full_html=True)


def main():
    n = 240
    x, U, target, err, energy, dt, h = upwind(n, record=True)
    fig = go.Figure()
    for k in range(2):
        fig.add_trace(go.Scatter(x=x, y=U[k], mode='lines',
                     name=f'numeryczne u{k+1}'))
        fig.add_trace(go.Scatter(x=x, y=target[k], mode='lines',
                     line=dict(dash='dash'), name=f'dokladne u{k+1}'))
    fig.update_layout(title=f'Ukad symetryczny: mody o predkosciach -1 i 0.6; blad={err:.3g}',
                      xaxis_title='x', yaxis_title='u(t,x)')
    export(fig, 'pde21_profile_plotly.html')

    fig = go.Figure()
    fig.add_trace(go.Scatter(x=np.arange(energy.size) * dt, y=energy,
                             mode='lines', name='Dyskretna energia l2'))
    fig.update_layout(title='Kontrakcja energii: schemat pod prad, CFL=0.8',
                      xaxis_title='czas', yaxis_title='h sum |U_j|^2')
    export(fig, 'pde21_energia_plotly.html')
    assert np.max(np.diff(energy)) <= 1e-12, 'Energia powinna byc nierosnaca'
    assert np.allclose(A, A.T)

    sizes = np.array([40, 60, 90, 135, 200, 300, 450])
    errors = np.array([upwind(int(k))[3] for k in sizes])
    hs = 1.0 / sizes
    slope = np.polyfit(np.log(hs[-4:]), np.log(errors[-4:]), 1)[0]
    print('Rzad zbieznosci, cztery ostatnie siatki:', round(slope, 4))
    print('Blad l2 dla N=240:', round(err, 7))
    assert 0.65 < slope < 1.15, f'Nieoczekiwany rzad {slope}'
    fig = go.Figure()
    fig.add_trace(go.Scatter(x=hs, y=errors, mode='lines+markers', name='Blad l2'))
    fig.add_trace(go.Scatter(x=hs, y=errors[-1] * hs / hs[-1], mode='lines',
                             line=dict(dash='dash'), name='Skala O(h)'))
    fig.update_layout(title=f'Zbieznosc schematu pod prad: rzad {slope:.3f}',
                      xaxis_title='h', yaxis_title='blad l2',
                      xaxis_type='log', yaxis_type='log')
    export(fig, 'pde21_zbieznosc_plotly.html')


if __name__ == '__main__':
    main()
\end{lstlisting}

\section{Swobodna ewolucja Schrödingera na torusie — Plotly}
\label{sec:pde22-code}
Kod stosuje reprezentację spektralną na torusie, a nie swobodne
jądro propagatora na $\R^d$. Sprawdza zachowanie normy
$L^2$, ruch pakietu i wzór Parsevala na błąd obcięcia.

\begin{lstlisting}[style=pythonpde,caption={Ewolucja Schrodingera: pakiet, faza i błąd widmowy (Plotly)},label={code:pde22-plotly}]
"""Rozdzial 22: swobodny propagator Schrodingera na torusie.

Zaleznosci: numpy, plotly
Uruchomienie: python demo_schrodinger_plotly.py
Wyniki: pde22_pakiet_plotly.html, pde22_faza_plotly.html,
        pde22_obciecie_plotly.html
"""
from pathlib import Path
import numpy as np
import plotly.graph_objects as go

OUT = Path(__file__).resolve().parent
N = 512
x = -np.pi + (2 * np.pi / N) * np.arange(N)
k = np.arange(-100, 101)
k0, sigma, x0 = 12.0, 5.0, -1.0
coef = np.exp(-0.5 * ((k - k0) / sigma)**2) * np.exp(-1j * k * x0)

def psi(t, cutoff=None):
    mask = np.ones(k.shape, dtype=bool) if cutoff is None else np.abs(k) <= cutoff
    z = (coef[mask] * np.exp(-1j * t * k[mask]**2))[:, None]
    return np.sum(z * np.exp(1j * k[mask, None] * x[None, :]), axis=0) / np.sqrt(2 * np.pi)

def store(fig, name):
    fig.update_layout(template="plotly_white", font=dict(size=13),
                      margin=dict(l=70, r=30, t=70, b=60))
    fig.write_html(OUT / name, include_plotlyjs=True, full_html=True)

def main():
    times = [0.0, 0.015, 0.03, 0.06, 0.09]
    fig = go.Figure()
    for t in times:
        p = psi(t)
        fig.add_trace(go.Scatter(x=x, y=np.abs(p)**2, mode='lines',
                                 name=f't={t:.3f}'))
    fig.update_layout(title='Pakiet swobodny na torusie: propagacja i dyspersja',
                      xaxis_title='x', yaxis_title='|u(t,x)|^2')
    store(fig, 'pde22_pakiet_plotly.html')

    # Przestrzenna siatka nie zawiera prawego konca okresu:
    # okresowa kwadratura trapezowa jest dokladna dla pasmowych iloczynow.
    norms_periodic = [(2*np.pi/N)*np.sum(np.abs(psi(t))**2) for t in times]
    parseval = np.sum(np.abs(coef)**2)
    print('Maksymalna zmiana normy (kwadrat):',
          max(abs(v - parseval) for v in norms_periodic))
    assert max(abs(v-parseval) for v in norms_periodic) < 1e-10

    t = 0.04
    fig = go.Figure()
    fig.add_trace(go.Scatter(x=k, y=np.abs(coef), name='|c_k|', mode='lines'))
    fig.add_trace(go.Scatter(x=k, y=np.angle(np.exp(-1j*t*k*k)),
                             name='faza propagatora', yaxis='y2', mode='markers',
                             marker=dict(size=3)))
    fig.update_layout(title='Zachowany modul i zmieniajaca sie faza',
                      xaxis_title='k', yaxis_title='|c_k|',
                      yaxis2=dict(title='faza [rad]', overlaying='y', side='right'))
    store(fig, 'pde22_faza_plotly.html')

    full = psi(0.05)
    cutoffs = [6, 10, 14, 18, 24, 32, 48, 64]
    errors = [(2*np.pi/N*np.sum(np.abs(full-psi(0.05,K))**2))**0.5
              for K in cutoffs]
    exact_tails = [np.sqrt(np.sum(np.abs(coef[np.abs(k)>K])**2))
                   for K in cutoffs]
    assert np.allclose(errors, exact_tails, rtol=1e-7, atol=1e-12)
    fig = go.Figure()
    fig.add_trace(go.Scatter(x=cutoffs, y=errors, name='blad numeryczny',
                             mode='lines+markers'))
    fig.add_trace(go.Scatter(x=cutoffs, y=exact_tails, name='ogon Parsevala',
                             mode='lines', line=dict(dash='dash')))
    fig.update_layout(title='Obciecie szeregu i tozsamosc Parsevala',
                      xaxis_title='maksymalne |k|', yaxis_title='blad L2',
                      yaxis_type='log')
    store(fig, 'pde22_obciecie_plotly.html')

if __name__ == '__main__':
    main()
\end{lstlisting}



\section{Równanie Hamiltona--Jacobiego: charakterystyki i rozwiązania lepkościowe (Plotly)}
Kod oblicza wartości wzoru Hopfa--Laxa przez minimalizację na siatce,
rysuje charakterystyki i pochodną rozwiązania po powstaniu osobliwości.
Wykresy ilustrują teorię; nie zastępują dowodów z Rozdziału~\ref{chap:pde23}.
\begin{lstlisting}[style=pythonpde,caption={Wzór Hopfa--Laxa i powstawanie narożnika (Plotly)},label={code:pde23-plotly}]
from pathlib import Path
import numpy as np
import plotly.graph_objects as go
OUT = Path(__file__).resolve().parent
def hopf_lax(t: float, x: np.ndarray, y: np.ndarray):
    """Minimalizacja dla g(y)=cos(y) i H(p)=p^2/2."""
    cost = np.cos(y)[None, :] + (x[:, None] - y[None, :])**2 / (2 * t)
    argmin = np.argmin(cost, axis=1)
    return cost[np.arange(len(x)), argmin], y[argmin]
def save(fig, filename):
    fig.update_layout(template="plotly_white", font=dict(size=13))
    fig.write_html(OUT / filename, include_plotlyjs=True, full_html=True)
def main():
    x = np.linspace(-np.pi, np.pi, 601)
    y = np.linspace(-np.pi - 2.2, np.pi + 2.2, 1801)
    fig = go.Figure()
    for t in (0.2, 0.8, 1.02, 1.35):
        u, minimizer = hopf_lax(t, x, y)
        fig.add_scatter(x=x, y=u, mode="lines", name=f"t={t:.2f}")
    fig.add_scatter(x=x, y=np.cos(x), mode="lines", name="t=0", line=dict(dash="dash"))
    fig.update_layout(title="Rozwiazanie lepkosciowe: g(x)=cos x, H(p)=p^2/2",
                      xaxis_title="x", yaxis_title="u(t,x)")
    save(fig, "pde23_profile_plotly.html")
    yy = np.linspace(-np.pi, np.pi, 35)
    tt = np.linspace(0, 1.6, 125)
    fig = go.Figure()
    for a in yy:
        fig.add_scatter(x=a - tt * np.sin(a), y=tt, mode="lines",
                        showlegend=False, line=dict(width=1.2))
    fig.add_hline(y=1, line_dash="dash", annotation_text="pierwsza degeneracja: t=1")
    fig.update_layout(title="Charakterystyki: x(t,y)=y-t sin y",
                      xaxis_title="poozenie x", yaxis_title="czas t")
    save(fig, "pde23_charakterystyki_plotly.html")
    t = 1.35
    dx = x[1] - x[0]
    u, minimizer = hopf_lax(t, x, y)
    grad = np.gradient(u, dx)
    fig = go.Figure()
    fig.add_scatter(x=x, y=grad, mode="lines", name="gradient numeryczny")
    fig.add_scatter(x=x, y=(x - minimizer) / t, mode="lines",
                    name="(x-y_min)/t", line=dict(dash="dash"))
    fig.update_layout(title="Skok gradientu rozwiazania po t=1",
                      xaxis_title="x", yaxis_title="pochodna przestrzenna")
    save(fig, "pde23_skok_gradientu_plotly.html")
    k = np.argmin(np.abs(x))
    xx = x[k]
    costs = np.cos(y) + (xx - y)**2 / (2*t)
    near = np.where(costs <= costs.min() + 1e-4)[0]
    assert np.any(y[near] < -0.4) and np.any(y[near] > 0.4)
    assert np.all(1 - 0.8*np.cos(yy) > 0)
    print("Trzy wizualizacje zapisane. Test symetrycznych minimow: OK.")
if __name__ == "__main__":
    main()
\end{lstlisting}


\section{Skalarne prawa zachowania: rozwiązania Riemanna — Plotly}
\label{sec:pde24-code}
Program implementuje monotoniczny schemat Godunowa dla strumienia
$f(u)=u^2/2$, porównuje go ze ścisłymi rozwiązaniami
entropijnymi zagadnienia Riemanna i wyznacza błąd w normie
$L^1$ na skończonym przedziale obliczeniowym.
Dane Riemanna nie są całkowalne na całej prostej; eksperyment
wykonujemy w skończonym oknie przestrzennym.
\begin{lstlisting}[style=pythonpde,caption={Fala uderzeniowa, fala rozrzedzenia i metoda Godunowa — Plotly},label={code:pde24-plotly}]
"""Skalarne prawa zachowania: fala uderzeniowa, rozrzedzenie, zbieznosc.

Uruchomienie: python demo_prawa_zachowania_plotly.py
Wymagania: numpy, plotly
"""
from pathlib import Path
import numpy as np
import plotly.graph_objects as go

OUT = Path(__file__).resolve().parent


def flux(u):
    return 0.5 * np.asarray(u) ** 2


def godunov(left, right):
    """Strumien Godunowa dla wypukego f(u)=u^2/2."""
    left = np.asarray(left)
    right = np.asarray(right)
    rarefaction = np.where(left >= 0, flux(left),
                          np.where(right <= 0, flux(right), 0.0))
    shock = np.maximum(flux(left), flux(right))
    return np.where(left <= right, rarefaction, shock)


def riemann_exact(x, t, left, right):
    if left > right:
        speed = (left + right) / 2
        return np.where(x < speed*t, left, right)
    if t == 0:
        return np.where(x < 0, left, right)
    return np.where(x <= left*t, left,
                    np.where(x >= right*t, right, x/t))


def numerical(n, tfinal, left, right):
    dx = 6.0/n
    x = -3.0 + (np.arange(n) + 0.5) * dx
    u = np.where(x < 0, left, right).astype(float)
    t = 0.0
    while t < tfinal - 1e-13:
        max_speed = max(1.0, np.max(np.abs(u)))
        dt = min(0.43*dx/max_speed, tfinal-t)
        padded = np.concatenate(([left], u, [right]))
        numerical_flux = godunov(padded[:-1], padded[1:])
        u -= dt/dx * (numerical_flux[1:] - numerical_flux[:-1])
        t += dt
    return x, u, dx


def graph(filename, title, traces, xaxis='x', yaxis='u'):
    fig = go.Figure()
    for x, y, name, style in traces:
        fig.add_trace(go.Scatter(x=x, y=y, name=name, mode='lines',
                                 line=dict(dash=style)))
    fig.update_layout(template='plotly_white', title=title,
                      xaxis_title=xaxis, yaxis_title=yaxis)
    if filename.endswith('zbieznosc.html'):
        fig.update_xaxes(type='log')
        fig.update_yaxes(type='log')
    fig.write_html(OUT / filename, include_plotlyjs=True, full_html=True)


def main():
    t = 0.8
    for left, right, name in [(1.0, 0.0, 'uderzenie'),
                               (0.0, 1.0, 'rozrzedzenie')]:
        x, u, dx = numerical(600, t, left, right)
        exact = riemann_exact(x, t, left, right)
        graph(f'pde24_{name}.html', f'Burgers: {name} (t={t})',
              [(x, exact, 'Rozwiazanie entropijne', 'solid'),
               (x, u, 'Schemat Godunowa', 'dash')])
        print(name, 'L1 error:', np.sum(abs(u-exact))*dx)

    ns = np.array([100, 200, 400, 800, 1600])
    trace_data = []
    for left, right, name, dash in [(1.0, 0.0, 'Uderzenie', 'solid'),
                                   (0.0, 1.0, 'Rozrzedzenie', 'dash')]:
        errors = []
        for n in ns:
            x, u, dx = numerical(int(n), t, left, right)
            errors.append(np.sum(abs(u-riemann_exact(x,t,left,right)))*dx)
        errors = np.asarray(errors)
        rate = np.polyfit(np.log(6.0/ns), np.log(errors), 1)[0]
        trace_data.append((6.0/ns, errors, f'{name}, nachylenie {rate:.2f}', dash))
        print(f'{name}: measured rate {rate:.4f}')
    graph('pde24_zbieznosc.html', 'Schemat Godunowa: bad L1',
          trace_data, xaxis='Krok siatki h', yaxis='Bad L1')

    # Kontrola zachowania ograniczenia maksimum przez monotoniczny schemat.
    x, u, _ = numerical(800, t, 1.0, 0.0)
    assert np.all(u >= -1e-12) and np.all(u <= 1+1e-12)


if __name__ == '__main__':
    main()
\end{lstlisting}


\section{Półliniowe równania ewolucyjne: energia i wybuch — Plotly}
\label{sec:pde25-code}
Program ilustruje spadek energii w równaniu Allena--Cahna
oraz dolną barierę na średnią rozwiązania równania
$u_t=\Delta u+u^2$ na torusie. Schemat Stranga jest
eksperymentem numerycznym, a nie dowodem twierdzeń.
\begin{lstlisting}[style=pythonpde,caption={Równania półliniowe: energia i eksplozja (Plotly)},label={code:pde25-plotly}]
"""PDE, rozdzial 25: reakcja, energia i eksplozja na torusie.

Uruchomienie: python demo_pololiniowe_plotly.py
Zaleznosci: numpy, plotly
Zapisuje trzy samodzielne wykresy interaktywne do katalogu skryptu.
"""
from pathlib import Path
import numpy as np
import plotly.graph_objects as go

OUT = Path(__file__).resolve().parent
N = 512
x = 2 * np.pi * np.arange(N) / N
k = np.fft.fftfreq(N, d=1 / N)
dx = 2 * np.pi / N

def diffuse(u, dt):
    return np.fft.ifft(np.fft.fft(u) * np.exp(-(k * k) * dt)).real

def save(fig, name):
    fig.update_layout(template="plotly_white", font=dict(size=13))
    fig.write_html(OUT / name, include_plotlyjs=True, full_html=True)

def allen_cahn():
    # Splitting Stranga: dyfuzja, dokladny lokalny przeplyw u'=u-u^3,
    # dyfuzja. Dokladny przeplyw reakcji zachowuje znak.
    u = 0.17 + 0.72 * np.cos(x) + 0.14 * np.cos(3 * x)
    dt = 0.001
    snapshots = {0.0: u.copy()}
    ts, energies = [], []
    def energy(v):
        vx = np.fft.ifft(1j * k * np.fft.fft(v)).real
        return dx * np.sum(0.5 * vx**2 + 0.25 * (v*v - 1)**2)
    for j in range(1601):
        if j % 20 == 0:
            ts.append(j * dt)
            energies.append(energy(u))
        if j in (400, 800, 1600):
            snapshots[round(j * dt, 2)] = u.copy()
        if j == 1600:
            break
        u = diffuse(u, dt / 2)
        e2 = np.exp(2 * dt)
        u = e2**0.5 * u / np.sqrt(1 + u*u * (e2 - 1))
        u = diffuse(u, dt / 2)
    # Wybieramy dostatecznie dokladny krok, by energia byla nierosnaca.
    assert np.max(np.diff(energies)) < 1e-6
    fig = go.Figure()
    for t, v in snapshots.items():
        fig.add_trace(go.Scatter(x=x, y=v, name=f"t={t:g}"))
    fig.update_layout(title="Rownanie Allena-Cahna: wygadzanie i reakcja", xaxis_title="x", yaxis_title="u(t,x)")
    save(fig, "pde25_allen_cahn_profile.html")
    fig2 = go.Figure(go.Scatter(x=ts, y=energies, name="E(u(t))"))
    fig2.update_layout(title="Rownanie Allena-Cahna: dyssypacja energii", xaxis_title="czas t", yaxis_title="energia")
    save(fig2, "pde25_allen_cahn_energia.html")
    print("Rownanie Allena-Cahna: energia poczatkowa i koncowa:", energies[0], energies[-1])

def blowup_mean():
    # u_t=Delta u+u^2, u0=0.35+0.15 cos(x), dodatnia srednia 0.35.
    u = 0.35 + 0.15 * np.cos(x)
    m0 = np.mean(u)
    bound = 1.0 / m0
    dt = 0.0005
    tmax = 2.2
    ts, means, mins, maxs = [], [], [], []
    for j in range(int(round(tmax/dt)) + 1):
        if j % 40 == 0:
            ts.append(j * dt)
            means.append(float(np.mean(u)))
            mins.append(float(np.min(u)))
            maxs.append(float(np.max(u)))
        if j == int(round(tmax/dt)):
            break
        u = diffuse(u, dt / 2)
        assert dt * np.max(u) < 1
        u = u / (1 - dt * u)  # scisly przeplyw reakcji u'=u^2
        u = diffuse(u, dt / 2)
    ts = np.asarray(ts)
    barrier = m0 / (1 - m0 * ts)
    # Blad splittigu i kwadratury dopuszczamy ponizej 1e-3.
    assert np.min(np.asarray(means) - barrier) > -1e-3
    fig = go.Figure()
    fig.add_trace(go.Scatter(x=ts, y=means, name="srednia numeryczna"))
    fig.add_trace(go.Scatter(x=ts, y=barrier, name="bariera Jensena / ODE", line=dict(dash="dash")))
    fig.add_trace(go.Scatter(x=ts, y=maxs, name="maksimum"))
    fig.update_layout(title=f"u_t=Delta u + u^2: srednia i dolna bariera, T*  {bound:.3f}", xaxis_title="czas t", yaxis_title="wartosc")
    save(fig, "pde25_eksplozja_srednia.html")
    print("Bariera T* <=", bound, "roznica sredniej i bariery w koncu", means[-1] - barrier[-1])

if __name__ == "__main__":
    allen_cahn()
    blowup_mean()

\end{lstlisting}



\section{Nieliniowy strumień p-laplasjanu: przykłady — Plotly}
\label{sec:pde26-code}
Program rozwiązujący jednowymiarowe zagadnienie $-\Delta_pu=1$
z zerowymi warunkami Dirichleta konstruuje dyskretny strumień
z równania bilansu i odtwarza węzłowe wartości rozwiązania.
Osobno mierzymy błąd i wzrost przybliżonej drugiej pochodnej
w maksimum. Przy $p>2$ rozwiązanie jest klasy $C^1$,
lecz nie musi należeć do $C^2$. Obliczenia stanowią
ilustrację, a nie dowód twierdzenia Minty'ego--Browdera.
\begin{lstlisting}[style=pythonpde,caption={Nieliniowy strumień, błąd aproksymacji i krzywizna (Plotly)},label={code:pde26-plotly}]
"""PDE 26: p-Laplacian Dirichlet, finite-volume flux and Plotly.
Requires: numpy, plotly. Run: python demo_plaplasjan_plotly.py
"""
from pathlib import Path
import numpy as np
import plotly.graph_objects as go

OUT = Path(__file__).resolve().parent
PS = (1.5, 2.0, 3.0, 5.0)


def exact(x, p):
    exponent = p / (p - 1)
    return ((p - 1) / p) * (0.5**exponent - np.abs(x - 0.5)**exponent)


def discrete(p, n):
    """Solve -D_h(|D_h u|^(p-2) D_h u)=1 with zero end values.
    In 1D the numerical flux is q_{i+1/2}=1/2-(i+1/2)h.
    """
    x = np.linspace(0, 1, n + 1)
    h = 1.0 / n
    xm = (x[1:] + x[:-1]) / 2
    flux = 0.5 - xm
    grad = np.sign(flux) * np.abs(flux)**(1 / (p - 1))
    values = np.r_[0.0, np.cumsum(h * grad)]
    assert abs(values[-1]) < 5e-13
    q_recovered = np.sign(np.diff(values) / h) * np.abs(np.diff(values) / h)**(p - 1)
    residual = -(np.diff(q_recovered) / h) - 1
    assert np.max(np.abs(residual)) < 1e-8
    return x, values


def save(fig, filename):
    fig.update_layout(template='plotly_white', font=dict(size=13))
    fig.write_html(OUT / filename, include_plotlyjs=True, full_html=True)


def make_profiles():
    x_dense = np.linspace(0, 1, 1001)
    fig = go.Figure()
    for p in PS:
        xd, ud = discrete(p, 16)
        fig.add_trace(go.Scatter(x=x_dense, y=exact(x_dense, p),
                                 mode='lines', name=f'p={p:g}, exact'))
        fig.add_trace(go.Scatter(x=xd, y=ud, mode='markers',
                                 name=f'p={p:g}, grid n=16'))
    fig.update_layout(title='One-dimensional p-Laplacian: -Delta_p u = 1',
                      xaxis_title='x', yaxis_title='u(x)')
    save(fig, 'pde26_profile.html')


def make_convergence():
    ns = np.array((20, 40, 80, 160, 320, 640))
    fig = go.Figure()
    for p in (1.5, 3.0, 5.0):
        errors = []
        for n in ns:
            x, u = discrete(p, int(n))
            errors.append(np.max(np.abs(u - exact(x, p))))
        errors = np.asarray(errors)
        rate = -np.polyfit(np.log(ns), np.log(errors), 1)[0]
        print('p=', p, 'convergence rate=', float(rate))
        assert errors[-1] < errors[0]
        fig.add_trace(go.Scatter(x=1/ns, y=errors, mode='lines+markers',
                                 name=f'p={p:g}, slope={rate:.2f}'))
    fig.update_layout(title='Error of nonlinear flux discretization',
                      xaxis_title='grid step h', yaxis_title='max nodal error',
                      xaxis_type='log', yaxis_type='log')
    save(fig, 'pde26_zbieznosc.html')


def make_curvature():
    ns = np.array((20, 40, 80, 160, 320, 640))
    fig = go.Figure()
    for p in (2.0, 3.0, 5.0):
        alpha = p / (p - 1)
        c = (p - 1) / p
        hs = 1 / ns
        # Exact centered second difference of the continuum solution.
        curvature = 2 * c * hs**(alpha - 2)
        assert curvature[-1] >= curvature[0] - 1e-12
        fig.add_trace(go.Scatter(x=hs, y=curvature, mode='lines+markers',
                                 name=f'p={p:g}, power={alpha-2:.3f}'))
    fig.update_layout(title='Growth of the centered second difference at x=1/2',
                      xaxis_title='h', yaxis_title='-D_h^2 u(1/2)',
                      xaxis_type='log', yaxis_type='log')
    save(fig, 'pde26_krzywizna.html')


if __name__ == '__main__':
    make_profiles()
    make_convergence()
    make_curvature()
    print('Discrete flux residuals and endpoint conditions verified.')
\end{lstlisting}


\section{Reakcja--dyfuzja: fale biegnące, energia i niestabilność Turinga (Plotly)}
\label{sec:pde27-code}
Program symuluje fale biegnące dla równania Fishera--KPP,
rozpad energii w równaniu Allena--Cahna i wzrost modów
w niestabilności Turinga. Wizualizacje ilustrują wyniki analityczne;
nie zastępują ich dowodów.
\begin{lstlisting}[style=pythonpde,caption={Fale reakcji--dyfuzji, energia i niestabilność Turinga (Plotly)},label={code:pde27-plotly}]
"""PDE 27: Fisher--KPP, fala biegnaca, Allen--Cahn i Turing.

Uruchomienie: pip install numpy plotly; python demo_reakcja_dyfuzja_plotly.py
Wyniki: pde27_fronty.html, pde27_fala_jawna.html,
        pde27_energia.html, pde27_turing.html
"""
from pathlib import Path
import numpy as np
import plotly.graph_objects as go

OUT = Path(__file__).resolve().parent


def save(fig, filename):
    fig.update_layout(template="plotly_white", font=dict(size=12),
                      legend=dict(orientation="h", y=-0.22))
    fig.write_html(OUT / filename, include_plotlyjs=True, full_html=True)


def reaction_kpp(u, dt, r=1.0):
    e = np.exp(r * dt)
    return u * e / (1 + u * (e - 1))


def split_step_kpp():
    n, L, dt, tmax = 1024, 120.0, 0.04, 16.0
    x = np.linspace(-L/2, L/2, n, endpoint=False)
    k = 2 * np.pi * np.fft.fftfreq(n, d=L/n)
    diffusion_half = np.exp(-k*k*dt/2)
    u = .8 * np.exp(-x*x/6)
    fig = go.Figure()
    fig.add_scatter(x=x, y=u, name="t=0")
    for j in range(1, round(tmax/dt)+1):
        u = np.fft.ifft(np.fft.fft(u) * diffusion_half).real
        u = reaction_kpp(np.clip(u, 0, 1), dt)
        u = np.fft.ifft(np.fft.fft(u) * diffusion_half).real
        u = np.clip(u, 0, 1)
        if j % 100 == 0:
            fig.add_scatter(x=x, y=u, name=f"t={j*dt:g}")
    assert np.min(u) >= 0 and np.max(u) <= 1
    fig.update_layout(title="Rownanie Fishera-KPP: rozprzestrzenianie danych skupionych",
                      xaxis_title="x", yaxis_title="u(t,x)")
    save(fig, 'pde27_fronty.html')


def explicit_front():
    z = np.linspace(-28, 28, 1601)
    w = np.exp(z / np.sqrt(6))
    u = (1+w)**-2
    up = -2*w/(np.sqrt(6)*(1+w)**3)
    upp = -w/(3*(1+w)**3)+w*w/(1+w)**4
    c = 5/np.sqrt(6)
    residual = upp+c*up+u*(1-u)
    assert np.max(np.abs(residual)) < 5e-15
    fig = go.Figure()
    for t in (0, 3, 6, 9):
        x = z+c*t
        fig.add_scatter(x=x, y=u, name=f"t={t}")
    fig.update_layout(title=f"Fala rownania Fishera-KPP: c=5/sqrt6~{c:.4f}, residuum < 5e-15",
                      xaxis_title="x", yaxis_title="U(x-ct)")
    save(fig, 'pde27_fala_jawna.html')
    return float(np.max(np.abs(residual)))


def reaction_allen(u, dt):
    e = np.exp(dt)
    return u*e/np.sqrt(1+u*u*(e*e-1))


def allen_cahn_energy():
    n, L, D, dt = 512, 20.0, .3, .02
    x = np.linspace(0, L, n, endpoint=False)
    k = 2*np.pi*np.fft.fftfreq(n, d=L/n)
    half = np.exp(-D*k*k*dt/2)
    u = .6*np.cos(2*np.pi*x/L)+.15*np.cos(6*np.pi*x/L)
    energy, time = [], []
    for j in range(601):
        if j % 5 == 0:
            ux = np.fft.ifft(1j*k*np.fft.fft(u)).real
            E = (L/n)*np.sum(D*ux**2/2+(u*u-1)**2/4)
            energy.append(E)
            time.append(j*dt)
        if j < 600:
            u = np.fft.ifft(np.fft.fft(u)*half).real
            u = reaction_allen(u, dt)
            u = np.fft.ifft(np.fft.fft(u)*half).real
    energy = np.asarray(energy)
    assert energy[-1] < energy[0] and np.all(np.isfinite(energy))
    fig = go.Figure(go.Scatter(x=time, y=energy, mode="lines+markers", name="E(u)"))
    fig.update_layout(title="Rownanie Allena-Cahna: dyssypacja energii (rozszczepienie Stranga)",
                      xaxis_title="t", yaxis_title="energia")
    save(fig, 'pde27_energia.html')
    return float(np.max(np.diff(energy)))


def turing_mode():
    mu = np.linspace(0, 1.1, 400)
    lam = []
    for m in mu:
        A = np.array([[1-m,-2],[2,-3-10*m]])
        lam.append(np.max(np.linalg.eigvals(A).real))
    lam = np.asarray(lam)
    at = np.linalg.eigvals(np.array([[.7,-2],[2,-6]])).real
    assert max(at) > 0
    assert np.all(lam[(mu > .2) & (mu < .5)] > 0)
    fig = go.Figure(go.Scatter(x=mu, y=lam, name="max Re lambda(J-muD)"))
    fig.add_hline(y=0, line_dash="dash")
    fig.add_vrect(x0=.2, x1=.5, fillcolor="lightgray", opacity=.18,
                  line_width=0, annotation_text="niestabilne mody")
    fig.update_layout(title="Niestabilnosc Turinga: stabilna reakcja, niestabilna dyfuzja",
                      xaxis_title="mu (wartosc wasna -Delta)", yaxis_title="tempo wzrostu")
    save(fig, 'pde27_turing.html')
    return float(max(at))


if __name__ == '__main__':
    split_step_kpp()
    print('maks. residuum jawnej fali:', explicit_front())
    print('maks. przyrost energii:', allen_cahn_energy())
    print('dodatnia wartosc wlasna Turinga:', turing_mode())
    print('Zapisano 4 wykresy Plotly w:', OUT)

\end{lstlisting}


\section{Równania Naviera--Stokesa: wir Taylora--Greena i rzut Leraya — Plotly}
\label{sec:pde28-code}
Poniższy program sprawdza dokładną energię wiru Taylora--Greena,
przedstawia wirowość i ciśnienie oraz oblicza rzut Leraya
w przestrzeni Fouriera. W modelowym polu dywergencja po rzutowaniu
jest bliska zeru z dokładnością maszynową. To eksperyment numeryczny,
nie zamiennik dowodu istnienia rozwiązania Leraya--Hopfa.
\begin{lstlisting}[style=pythonpde,caption={Wir Taylora--Greena i rzut Leraya (Plotly)},label={code:pde28-plotly}]
"""Rownania Naviera-Stokesa: wir Taylora-Greena i rzut Leraya.

pip install numpy plotly
python demo_navier_stokes_plotly.py
Wyniki: pde28_wiry.html, pde28_energia.html,
        pde28_cisnienie.html, pde28_rzut_helmholtza.html
"""
from pathlib import Path
import numpy as np
import plotly.graph_objects as go

OUT = Path(__file__).resolve().parent
N, NU = 128, 0.08
x = np.linspace(0, 2*np.pi, N, endpoint=False)
X, Y = np.meshgrid(x, x, indexing='ij')
k = np.fft.fftfreq(N, d=1/N)  # Fourier integers, not physical cycles
KX, KY = np.meshgrid(k, k, indexing='ij')


def store(fig, name):
    fig.update_layout(template='plotly_white', font=dict(size=13))
    fig.write_html(OUT / name, include_plotlyjs=True, full_html=True)


def taylor_green(t):
    decay = np.exp(-2*NU*t)
    ux = decay*np.sin(X)*np.cos(Y)
    uy = -decay*np.cos(X)*np.sin(Y)
    p = 0.25*decay**2*(np.cos(2*X)+np.cos(2*Y))
    omega = 2*decay*np.sin(X)*np.sin(Y)
    return ux, uy, p, omega


def divergence_fft(vx, vy):
    return np.fft.ifft2(1j*KX*np.fft.fft2(vx)
                          +1j*KY*np.fft.fft2(vy)).real


def project(vx, vy):
    hx, hy = np.fft.fft2(vx), np.fft.fft2(vy)
    den = KX*KX+KY*KY
    safe = np.where(den == 0, 1, den)
    dot = KX*hx+KY*hy
    hx = hx-KX*dot/safe
    hy = hy-KY*dot/safe
    hx[0,0]=hy[0,0]=0
    return np.fft.ifft2(hx).real, np.fft.ifft2(hy).real


def plots():
    u1,u2,p,omega = taylor_green(0.5)
    fig = go.Figure(go.Contour(x=x,y=x,z=omega.T, colorscale='RdBu',
                 contours=dict(showlabels=True)))
    fig.update_layout(title='Wir Taylora-Greena: wirowosc w chwili t=0.5',
                      xaxis_title='x',yaxis_title='y', yaxis_scaleanchor='x')
    store(fig, 'pde28_wiry.html')

    times=np.linspace(0,10,150)
    e_exact=np.pi**2*np.exp(-4*NU*times)
    e_numeric=np.array([0.5*(taylor_green(t)[0]**2+
                           taylor_green(t)[1]**2).mean()*(2*np.pi)**2
                        for t in times])
    assert np.max(np.abs(e_exact-e_numeric))<2e-12
    fig=go.Figure()
    fig.add_trace(go.Scatter(x=times,y=e_numeric,name='Energia z siatki',mode='markers',
                             marker=dict(size=4)))
    fig.add_trace(go.Scatter(x=times,y=e_exact,name='Energia dokladna',mode='lines'))
    fig.update_layout(title='Tozsamosc energii: E(t)=pi^2 exp(-4 nu t)',
                      xaxis_title='czas t',yaxis_title='E(t)')
    store(fig,'pde28_energia.html')

    fig=go.Figure(go.Heatmap(x=x,y=x,z=p.T,colorscale='RdBu'))
    fig.update_layout(title='Cisnienie wiru Taylora-Greena w chwili t=0.5',
                      xaxis_title='x',yaxis_title='y',yaxis_scaleanchor='x')
    store(fig,'pde28_cisnienie.html')

    fx=np.sin(Y)+np.cos(X)
    fy=np.sin(X)+np.cos(Y)
    px,py=project(fx,fy)
    assert np.max(np.abs(px-np.sin(Y)))<2e-12
    assert np.max(np.abs(py-np.sin(X)))<2e-12
    div_before=np.max(np.abs(divergence_fft(fx,fy)))
    div_after=np.max(np.abs(divergence_fft(px,py)))
    assert div_before>1.9 and div_after<1e-10
    assert np.max(np.abs(divergence_fft(u1,u2)))<1e-10
    fig=go.Figure()
    step=9
    for vx,vy,name,dx in [(fx,fy,'F przed rzutem',0), (px,py,'PF po rzucie',2*np.pi+1)]:
        first=True
        for i in range(0,N,step):
            for j in range(0,N,step):
                xi,yi=x[i]+dx,x[j]
                fig.add_trace(go.Scatter(x=[xi,xi+0.27*vx[i,j],None],
                      y=[yi,yi+0.27*vy[i,j],None],mode='lines',
                      line=dict(width=1.1),name=name,showlegend=first))
                first=False
    fig.update_layout(title=f'Rzut Leraya: max div przed={div_before:.2f}, po={div_after:.1e}',
                      xaxis_title='dwa pola obok siebie',yaxis_title='y',
                      yaxis_scaleanchor='x')
    store(fig,'pde28_rzut_helmholtza.html')
    print('Energia: maksymalny blad', np.max(np.abs(e_exact-e_numeric)))
    print('Dywergencja przed i po rzucie',div_before,div_after)


if __name__ == '__main__':
    plots()
\end{lstlisting}


\section{Zagadnienie przeszkody i reakcja kontaktowa — Plotly}
\label{sec:pde29-code}
Program implementuje metodę Gaussa--Seidela z projekcją
na zbiór funkcji spełniających przeszkodę. Porównuje wyniki
z dokładnym rozwiązaniem z przykładu~\ref{ex:pde29-explicit}
i kontroluje nieujemność reakcji oraz warunek komplementarności.
Listing jest obramowany; samodzielny plik \texttt{demo\_przeszkoda\_plotly.py}
znajduje się również w pakiecie projektu.
\begin{lstlisting}[style=pythonpde,caption={Przeszkoda, reakcja i zbieżność schematu (Plotly)},label={code:pde29-plotly}]
"""Rozdzial 29 PDE: metoda Gaussa-Seidela z projekcja dla przeszkody.

Wymagania: pip install numpy plotly
Wykonanie: python demo_przeszkoda_plotly.py
Wyniki: pde29_przeszkoda.html, pde29_reakcja.html,
        pde29_zbieznosc.html
"""
from pathlib import Path
import numpy as np
import plotly.graph_objects as go

OUT = Path(__file__).resolve().parent


def obstacle(x: np.ndarray) -> np.ndarray:
    return (1.0 - x * x) ** 2


def exact(x: np.ndarray) -> np.ndarray:
    a = np.abs(x)
    return np.where(a <= 1 / 3, (1.0 - x * x) ** 2,
                    32 / 27 * (1.0 - a))


def projected_gauss_seidel(n: int, omega=1.75, tol=2e-12,
                            max_sweeps=18000):
    h = 2.0 / n
    x = np.linspace(-1.0, 1.0, n + 1)
    psi = obstacle(x)
    u = psi.copy()
    u[0] = u[-1] = 0
    history = []
    for k in range(max_sweeps):
        change = 0.0
        for i in range(1, n):
            new = max(psi[i], u[i] + omega *
                      (0.5 * (u[i - 1] + u[i + 1]) - u[i]))
            change = max(change, abs(new - u[i]))
            u[i] = new
        if k % 25 == 0:
            energy = np.sum(np.diff(u) ** 2) / (2 * h)
            history.append((k, energy, change))
        if change < tol:
            break
    else:
        raise RuntimeError(f'Projected GS did not converge for n={n}')
    reaction = (2 * u[1:-1] - u[:-2] - u[2:]) / (h * h)
    slack = u[1:-1] - psi[1:-1]
    return x, u, psi, reaction, slack, np.array(history), k + 1


def save(fig, name):
    fig.update_layout(template='plotly_white', font=dict(size=13))
    fig.write_html(OUT / name, include_plotlyjs=True, full_html=True)


def main():
    ns = [40, 80, 160]
    results = [projected_gauss_seidel(n) for n in ns]
    errors = []
    for n, (x, u, psi, reaction, slack, hist, sweeps) in zip(ns, results):
        err = np.max(np.abs(u - exact(x)))
        errors.append(err)
        assert np.min(slack) > -1e-11, (n, np.min(slack))
        assert np.min(reaction) > -2e-7, (n, np.min(reaction))
        assert np.max(np.abs(reaction * slack)) < 5e-7
        print(f'N={n}, sweeps={sweeps}, max error={err:.4e}, '
              f'min reaction={reaction.min():.2e}, '
              f'max complementarity={np.max(np.abs(reaction * slack)):.2e}')

    x, u, psi, reaction, slack, hist, _ = results[-1]
    xf = np.linspace(-1, 1, 901)
    fig = go.Figure()
    fig.add_trace(go.Scatter(x=xf, y=obstacle(xf),
                             name='Przeszkoda psi', mode='lines'))
    fig.add_trace(go.Scatter(x=xf, y=exact(xf),
                             name='Rozwi\u0105zanie dok\u0142adne', mode='lines'))
    fig.add_trace(go.Scatter(x=x, y=u,
                             name='Wynik numeryczny', mode='markers',
                             marker=dict(size=5)))
    for point in (-1 / 3, 1 / 3):
        fig.add_vline(x=point, line_dash='dash', line_color='gray')
    fig.update_layout(title='Zagadnienie przeszkody: kontakt i swobodny brzeg',
                      xaxis_title='x', yaxis_title='Wysoko\u015b\u0107')
    save(fig, 'pde29_przeszkoda.html')

    xnodes = x[1:-1]
    reaction_exact = np.where(np.abs(xnodes) < 1 / 3,
                              4 - 12 * xnodes * xnodes, 0.0)
    fig = go.Figure()
    fig.add_trace(go.Scatter(x=xnodes, y=reaction,
                             mode='lines+markers',
                             name='Reakcja dyskretna',
                             marker=dict(size=4)))
    fig.add_trace(go.Scatter(x=xnodes, y=reaction_exact,
                             name='Reakcja dok\u0142adna', mode='lines'))
    fig.update_layout(title='Reakcja przeszkody: -u" >= 0',
                      xaxis_title='x', yaxis_title='Mno\u017cnik ograniczenia')
    save(fig, 'pde29_reakcja.html')

    fig = go.Figure()
    fig.add_trace(go.Scatter(x=np.array(ns), y=errors, mode='lines+markers',
                             name='Maksymalny b\u0142\u0105d w\u0119z\u0142owy'))
    fig.update_layout(title='Zbie\u017cno\u015b\u0107 metody z projekcj\u0105',
                      xaxis_type='log', yaxis_type='log',
                      xaxis_title='Liczba przedzia\u0142\xf3w N',
                      yaxis_title='B\u0142\u0105d maksimum')
    save(fig, 'pde29_zbieznosc.html')
    if all(err > 0 for err in errors):
        orders = np.log(np.array(errors[:-1]) / np.array(errors[1:])) / np.log(2)
        print('Obserwowane rz\u0119dy zbie\u017cno\u015bci:', orders)


if __name__ == '__main__':
    main()
\end{lstlisting}

\section{Homogenizacja periodyczna: rozwiązanie i strumień — Plotly}
\label{sec:pde30-code}
Program porównuje dokładne rozwiązania jednowymiarowe z równaniem
uśrednionym, pokazuje oscylacje gradientu oraz zbieżność błędu.
Ostatni wykres przedstawia anizotropię macierzy efektywnej
materiału dwuwarstwowego; jest ilustracją prawa
~\eqref{eq:pde30-layers}, a nie obliczeniem metodą elementów skończonych.
\begin{lstlisting}[style=pythonpde,caption={Homogenizacja okresowa i macierz efektywna (Plotly)},label={code:pde30-plotly}]
"""PDE, rozdzial 30: 1D homogenizacja okresowa i srednie efektywne.

Wymagania: pip install numpy plotly
Uruchomienie: python demo_homogenizacja_plotly.py
Pliki wynikowe: pde30_profile.html, pde30_gradient_strumien.html,
                pde30_zbieznosc.html, pde30_anizotropia.html
"""
from pathlib import Path
import numpy as np
import plotly.graph_objects as go

DIR = Path(__file__).resolve().parent
A_HOM = np.sqrt(3.0)


def cumulative_trapezoid(v, h):
    return np.concatenate(([0.0], np.cumsum((v[1:] + v[:-1]) * (h / 2.0))))


def micro_solution(periods, samples_per_period=240):
    eps = 1.0 / periods
    x = np.linspace(0.0, 1.0, periods * samples_per_period + 1)
    h = x[1] - x[0]
    a = 2.0 + np.sin(2.0 * np.pi * x / eps)
    b = 1.0 / a
    ib = cumulative_trapezoid(b, h)
    ibx = cumulative_trapezoid(b * x, h)
    c = ibx[-1] / ib[-1]
    u = c * ib - ibx
    gradient = b * (c - x)
    flux = c - x
    u0 = x * (1.0 - x) / (2.0 * A_HOM)
    grad0 = (0.5 - x) / A_HOM
    err = np.sqrt(np.trapezoid((u - u0) ** 2, x))
    return x, u, gradient, flux, u0, grad0, err


def export(fig, name):
    fig.update_layout(template="plotly_white", font=dict(size=13))
    fig.write_html(DIR / name, full_html=True, include_plotlyjs=True)


def main():
    first = go.Figure()
    for n in [4, 8, 16, 32]:
        x, u, du, q, u0, du0, err = micro_solution(n)
        skip = max(1, len(x) // 2000)
        first.add_trace(go.Scatter(x=x[::skip], y=u[::skip], name=f"eps=1/{n}", mode="lines"))
    first.add_trace(go.Scatter(x=x, y=u0, name="rozwiazanie graniczne", line=dict(color="black", dash="dash")))
    first.update_layout(title="Homogenizacja: rozwiazania mikroskopowe i graniczne",
                        xaxis_title="x", yaxis_title="u(x)")
    export(first, "pde30_profile.html")

    x, u, grad, flux, u0, grad0, err = micro_solution(24)
    fig = go.Figure()
    fig.add_trace(go.Scatter(x=x, y=grad, name="gradient u'eps"))
    fig.add_trace(go.Scatter(x=x, y=grad0, name="gradient u'0", line=dict(dash="dot")))
    fig.add_trace(go.Scatter(x=x, y=flux, name="strumien aepsu'eps", line=dict(width=3)))
    fig.add_trace(go.Scatter(x=x, y=0.5-x, name="strumien graniczny", line=dict(dash="dash")))
    fig.update_layout(title="Silna zbieznosc strumienia, oscylacje gradientu",
                      xaxis_title="x", yaxis_title="wartosc")
    export(fig, "pde30_gradient_strumien.html")

    ns = np.array([4, 8, 16, 32, 64], dtype=int)
    errors = np.array([micro_solution(int(n))[6] for n in ns])
    order = float(np.polyfit(np.log(1.0/ns), np.log(errors), 1)[0])
    fig = go.Figure()
    fig.add_trace(go.Scatter(x=1.0/ns, y=errors, mode="lines+markers", name="blad L^2"))
    fig.add_trace(go.Scatter(x=1.0/ns, y=errors[0]*(ns[0]/ns),
                             mode="lines", name="wzorzec O(eps)", line=dict(dash="dash")))
    fig.update_layout(title=f"Zbieznosc do rozwiazania granicznego; rzad ~{order:.3f}",
                      xaxis_type="log", yaxis_type="log",
                      xaxis_title="eps", yaxis_title="||ueps-u0|| L^2")
    export(fig, "pde30_zbieznosc.html")

    t = np.linspace(0, 2*np.pi, 360)
    fig = go.Figure()
    fig.add_trace(go.Scatter(x=np.sqrt(8.0/5.0)*np.cos(t),
                             y=np.sqrt(5.0/2.0)*np.sin(t),
                             name="warstwy: diag(8/5,5/2)"))
    fig.add_trace(go.Scatter(x=np.sqrt(2.5)*np.cos(t),
                             y=np.sqrt(2.5)*np.sin(t),
                             name="bledna srednia arytmetyczna", line=dict(dash="dash")))
    fig.update_yaxes(scaleanchor="x", scaleratio=1)
    fig.update_layout(title="Izotropowe fazy -> anizotropowa macierz efektywna",
                      xaxis_title="kierunek prostopadly do warstw (y1)",
                      yaxis_title="kierunek wzdluz warstw (y2)")
    export(fig, "pde30_anizotropia.html")
    assert abs(A_HOM - 1.7320508075688772) < 1e-12
    assert all(np.diff(errors) < 0), errors
    assert 0.8 < order < 1.2, order
    assert abs(flux[-1] - flux[0] + 1.0) < 1e-12
    assert abs(u[-1]) < 5e-12
    print("Efektywny wspolczynnik:", A_HOM)
    print("Rzad zbieznosci L2:", order)
    print("Bledy L2:", errors)


if __name__ == "__main__":
    main()
\end{lstlisting}

\section{Analiza mikrolokalna i symbole: eksperymenty Plotly}
\label{sec:pde31-code}
Rozkład częstości dla wąskiej warstwy, stożek charakterystyczny i parametryks.
\begin{lstlisting}[style=pythonpde,caption={Eksperymenty rozdziału 31 (Python i Plotly)},label={code:pde31-plotly}]
"""Rozdzial 31: kierunki osobliwosci i symbole Fourierowskie.
Wymaga numpy, plotly; zapisuje trzy interaktywne wykresy HTML.
Przyblizenie delta(x1) jest gladkie dla skonczonej szerokosci;
wykres ilustruje, a nie dowodzi twierdzenia o zbiorze falowym.
"""
from pathlib import Path
import numpy as np
import plotly.graph_objects as go

OUT = Path(__file__).resolve().parent


def save(fig, filename):
    fig.update_layout(template="plotly_white", font=dict(size=13))
    fig.write_html(OUT / filename, include_plotlyjs=True, full_html=True)


def directional_singularity():
    n = 256
    x = np.linspace(-np.pi, np.pi, n, endpoint=False)
    dx = x[1] - x[0]
    xx, yy = np.meshgrid(x, x, indexing="ij")
    angle = np.pi / 5
    normal = xx * np.cos(angle) + yy * np.sin(angle)
    # Okresowy profil waskiego pasa; zwezone przyblizenie delty.
    profile = np.exp(-0.5 * (np.sin(normal / 2) / 0.075) ** 2)
    window = np.exp(-0.5 * ((xx / 2.2) ** 8 + (yy / 2.2) ** 8))
    img = profile * window
    freq = np.fft.fftshift(np.fft.fftfreq(n, d=dx)) * 2 * np.pi
    amp = np.log1p(np.abs(np.fft.fftshift(np.fft.fft2(img))))
    fig = go.Figure(go.Heatmap(x=freq, y=freq, z=amp.T, colorscale="Viridis"))
    fig.update_layout(title="Kierunki wysokich czestotliwosci: waski pas",
                      xaxis_title="xi_1", yaxis_title="xi_2")
    fig.update_xaxes(range=[-45, 45]); fig.update_yaxes(range=[-45, 45])
    save(fig, "pde31_kierunki_fouriera.html")
    assert np.isfinite(amp).all()


def characteristic_cone():
    xi = np.linspace(-5, 5, 801)
    fig = go.Figure()
    fig.add_trace(go.Scatter(x=xi, y=2 * np.abs(xi), name="tau = +2|xi|"))
    fig.add_trace(go.Scatter(x=xi, y=-2 * np.abs(xi), name="tau = -2|xi|"))
    fig.update_layout(title="Zbior charakterystyczny rownania falowego, c=2",
                      xaxis_title="xi", yaxis_title="tau")
    save(fig, "pde31_charakterystyki.html")
    assert np.max(np.abs((2 * np.abs(xi)) ** 2 - 4 * xi ** 2)) < 1e-10


def parametrix_filter():
    xi = np.linspace(-20, 20, 1201)
    p = 1 + xi ** 2
    inverse = 1 / p
    fig = go.Figure()
    fig.add_trace(go.Scatter(x=xi, y=p, name="P(i xi)=1+xi^2"))
    fig.add_trace(go.Scatter(x=xi, y=inverse, name="q(xi)=1/(1+xi^2)"))
    fig.update_layout(title="Symbol eliptyczny i odwrotny mnoznik",
                      xaxis_title="xi", yaxis_title="wartosc")
    save(fig, "pde31_parametryks.html")
    assert np.allclose(p * inverse, 1)


if __name__ == "__main__":
    directional_singularity()
    characteristic_cone()
    parametrix_filter()
    print("Mikrolokalna: zapisano 3 wykresy Plotly.")
\end{lstlisting}


\section{Problemy odwrotne i sterowanie optymalne: eksperymenty Plotly}
\label{sec:pde32-code}
Zanik wartości singularnych, regularyzacja Tichonowa i sterowanie.
\begin{lstlisting}[style=pythonpde,caption={Eksperymenty rozdziału 32 (Python i Plotly)},label={code:pde32-plotly}]
"""Rozdzial 32: niestabilnosc, regularyzacja Tichonowa i sterowanie PDE.
Wymaga numpy, plotly. Zawiera testy formul dla modelu spektralnego.
"""
from pathlib import Path
import numpy as np
import plotly.graph_objects as go

OUT = Path(__file__).resolve().parent


def save(fig, filename):
    fig.update_layout(template="plotly_white", font=dict(size=13))
    fig.write_html(OUT / filename, include_plotlyjs=True, full_html=True)


def singular_values():
    n = np.arange(1, 41)
    sig_ell = 1 / (1 + n ** 2)
    sig_heat = np.exp(-0.01 * n ** 2)
    fig = go.Figure()
    fig.add_trace(go.Scatter(x=n, y=sig_ell, mode="lines+markers", name="(-Delta+I)^(-1)"))
    fig.add_trace(go.Scatter(x=n, y=sig_heat, mode="lines+markers", name="exp(T Delta), T=0.01"))
    fig.update_layout(title="Male wartosci singularne dwoch operatorow PDE",
                      xaxis_title="numer modu", yaxis_title="wartosc singularna",
                      yaxis_type="log")
    save(fig, "pde32_wartosci_singularne.html")


def inverse_heat():
    rng = np.random.default_rng(713)
    n = np.arange(1, 51)
    T = 0.007
    sigma = np.exp(-T * n ** 2)
    truth = np.exp(-n / 7) * np.sin(n / 3)
    delta = 0.002
    noise = rng.normal(0, 1, n.size)
    noise = delta * noise / np.linalg.norm(noise)
    data = sigma * truth + noise
    alpha = np.logspace(-10, -2, 130)
    recon = sigma[:, None] * data[:, None] / (sigma[:, None] ** 2 + alpha[None, :])
    error = np.linalg.norm(recon - truth[:, None], axis=0)
    idx = int(np.argmin(error)); alpha_best = alpha[idx]
    fig = go.Figure()
    fig.add_trace(go.Scatter(x=alpha, y=error, name="blad rekonstrukcji"))
    fig.add_vline(x=alpha_best, line_dash="dash",
                  annotation_text="minimum dla danych testowych")
    fig.update_layout(title="Regularyzacja Tichonowa: kompromis szum-gladzenie",
                      xaxis_title="alpha", xaxis_type="log",
                      yaxis_title="blad wspolczynnikow")
    save(fig, "pde32_regularyzacja.html")
    fig2 = go.Figure()
    fig2.add_trace(go.Scatter(x=n, y=truth, name="sygnal prawdziwy"))
    fig2.add_trace(go.Scatter(x=n, y=recon[:, idx], name="rekonstrukcja"))
    fig2.update_layout(title=f"Odtwarzanie modow, alpha={alpha_best:.2e}",
                       xaxis_title="numer modu", yaxis_title="wspolczynnik")
    save(fig2, "pde32_rekonstrukcja.html")
    assert error[idx] <= error[0] and error[idx] <= error[-1]
    return float(alpha_best), float(error[idx])


def elliptic_control():
    x = np.linspace(0, np.pi, 601)
    alpha = 0.15
    n = np.array([1.0, 3.0])
    desired_coeff = np.array([1.0, 0.35])
    sigma = 1 / (n ** 2)
    control_coeff = sigma * desired_coeff / (sigma ** 2 + alpha)
    observed_coeff = sigma * control_coeff
    sin_modes = np.sin(n[:, None] * x)
    y_desired = desired_coeff @ sin_modes
    y_state = observed_coeff @ sin_modes
    control = control_coeff @ sin_modes
    fig = go.Figure()
    for values, name in ((y_desired, "stan pozadany"),
                         (y_state, "stan optymalny"), (control, "sterowanie")):
        fig.add_trace(go.Scatter(x=x, y=values, name=name))
    fig.update_layout(title="Sterowanie rozlozone dla -y''=u",
                      xaxis_title="x", yaxis_title="wartosc")
    save(fig, "pde32_sterowanie.html")
    assert np.allclose((sigma ** 2 + alpha) * control_coeff,
                       sigma * desired_coeff)


if __name__ == "__main__":
    singular_values()
    a, err = inverse_heat()
    elliptic_control()
    print(f"Optymalne testowe alpha={a:.3e}, blad={err:.4f}; zapisano 4 wykresy.")
\end{lstlisting}


